\documentclass[11pt]{article}
\ifdefined\XeTeXversion
\else
  \pdfoutput=1
\fi

\usepackage[T1]{fontenc}
\usepackage{lmodern}
\usepackage[margin=1in]{geometry}
\usepackage{amsmath,amssymb,amsthm,mathtools,bm}
\usepackage{booktabs}
\usepackage{enumitem}
\usepackage[numbers,sort&compress]{natbib}
\usepackage{xcolor}
\usepackage{microtype}
\usepackage[colorlinks=true,allcolors=blue!55!black]{hyperref}

\newtheorem{theorem}{Theorem}[section]
\newtheorem{proposition}[theorem]{Proposition}
\newtheorem{corollary}[theorem]{Corollary}
\newtheorem{lemma}[theorem]{Lemma}
\theoremstyle{definition}
\newtheorem{definition}[theorem]{Definition}
\newtheorem{assumption}[theorem]{Assumption}
\theoremstyle{remark}
\newtheorem{remark}[theorem]{Remark}
\newtheorem{example}[theorem]{Example}

\newcommand{\simplex}{\Delta_{N}^{\circ}}
\newcommand{\R}{\mathbb{R}}
\newcommand{\E}{\mathbb{E}}
\newcommand{\KL}{D_{\mathrm{KL}}}
\newcommand{\LC}{\mathrm{LC}}
\newcommand{\FR}{\mathrm{FR}}
\newcommand{\AC}{\mathrm{AC}}
\newcommand{\op}{\mathrm{op}}
\newcommand{\Var}{\operatorname{Var}}
\newcommand{\Exp}{\operatorname{Exp}}
\newcommand{\Log}{\operatorname{Log}}
\newcommand{\dd}{\mathrm{d}}

\title{From Predictive Agreement to Geometric Agreement:\\
Stability of Fisher Connections and Semantic Transport in Neural Representations}
\author{Guy Basson}
\date{August 2026}

\begin{document}

\maketitle

\begin{abstract}
A probabilistic decoder whose predictive fiber quotient is a smooth manifold
equips that quotient with the pullback Fisher metric and Amari--Chentsov tensor.
We establish a stability hierarchy for comparing this geometry across models.
Exact equality of predictive images preserves every Amari $\alpha$-connection,
transport, curvature, and intrinsic predictive field. Uniform KL convergence
alone does not imply even pointwise Fisher-metric convergence. General
$\alpha$-transport is stable under uniform interiority, nondegeneracy, and
$C^2$ predictive agreement, but the resulting worst-case constants are poorly
adapted to large vocabularies. For the Levi--Civita connection we remove both
the vocabulary-size and minimum-token-probability dependence: $C^2$ closeness
of the square-root predictive maps controls the connection and transport, and
integrated KL plus a uniform Hilbert-valued Sobolev bound supplies this
closeness at an explicit rate. A paired boundary Bernoulli construction proves
that the same square-root regularity cannot control any fixed nonzero-$\alpha$
connection without additional information. Score-moment identities identify
the required raised third-order moment and yield probability-floor-free
formulas at minimal affine softmax heads. Finally, conditional Fisher variance
is the exact obstruction to descending a semantic effect through a
noninjective predictive map; it vanishes identically when that map is
injective. Constant vector reuse at an affine head, interpreted modulo its
predictive null space, is exponential transport on the ambient intervention
family, so superiority of any other connection is a held-out empirical
hypothesis rather than a consequence of pretraining.
\end{abstract}

\medskip
\noindent\textbf{MSC 2020:} 53B12, 62B11, 68T07.

\smallskip
\noindent\textbf{Keywords:} information geometry, Fisher--Rao metric, Amari $\alpha$-connections, parallel transport, stability estimates, neural representations.

\section{Introduction}

Modern neural networks compute with vectors while producing probability distributions. For an autoregressive language model, a hidden state $h$ entering the output head determines
\begin{equation}
    F(h)=p(\cdot\mid h)\in\simplex,
\end{equation}
where $\simplex$ is the interior of the categorical simplex over a vocabulary of size $N$. The hidden coordinate remains a point of $\R^d$. Nevertheless, its predictively visible component parameterizes a statistical model and inherits the Fisher geometry of the decoder.

Exact equality of identifiable predictive images gives an isomorphism of their induced statistical structures. Approximate equality at the loss level is weaker: cross-entropy controls probability values on the data distribution, whereas Fisher metrics and affine connections depend on derivatives of the predictive map. Geometric equivalence also does not imply that a semantic operation is determined by the next-token distribution, because hidden states can contain decoder-null or semantically relevant fiber directions.

At a linear softmax head, constant hidden-vector reuse is exactly the affine operation selected by the flat exponential connection. Ordinary addition is therefore one connection-dependent rule rather than an alternative to geometry. The resulting question is which connection, if any, makes a specified semantic effect transferable across contexts and models.

The analysis proceeds through six results:

\begin{enumerate}[leftmargin=*,itemsep=0.35em]
    \item Exact predictive naturality: two identifiable parameterizations of the same predictive family preserve the Fisher metric, Amari-Chentsov tensor, every $\alpha$-connection, parallel transport, and curvature up to reparameterization; every common smooth field on the shared predictive image pulls back equivariantly.
    \item Separation of predictive and geometric convergence: a smooth, uniformly regular Bernoulli family converges uniformly in KL while its Fisher metric fails to converge at a fixed point.
    \item Quantitative $C^k$ stability: on compact uniformly interior and uniformly identifiable families, $C^2$ predictive closeness controls the defect between aligned $\alpha$-connections and the commutator of their parallel transports, while $C^3$ closeness controls curvature.
    \item A sharp connection-selective boundary result: square-root regularity removes explicit vocabulary-size and minimum-probability factors from Levi--Civita bounds, while paired Bernoulli families converge in square-root $C^2$ even as every fixed nonzero-$\alpha$ connection defect diverges.
    \item Loss-to-geometry implications under explicit regularity: integrated KL plus uniform Sobolev control of either probability coordinates or square-root coordinates yields transport convergence; the latter result is vocabulary independent for Levi--Civita transport.
    \item Semantic conditional-variance decompositions: two orthogonal identities separately isolate failure to descend through a noninjective predictive map from connection mismatch and from cross-model conditional-mean-field mismatch.
\end{enumerate}

These results define a concrete comparison protocol: choose a predictive quotient and a cross-model alignment, estimate infinitesimal semantic effects on a declared continuous chart, and compare exponential, Levi--Civita, mixture, and learned-$\alpha$ transport through commutation error and held-out behavioral transfer. The square-root sphere representation itself is classical; the contribution here is the quantitative implication--separation package and its decoder interpretation.

\section{Related work}

Information geometry equips regular statistical models with the Fisher metric, the Amari--Chentsov cubic tensor, and dual affine connections \citep{amariNagaoka2000,amari2016,ayEtAl2018parametrized}. Invariance under sufficient statistics and the uniqueness of the Fisher and Amari--Chentsov tensors are classical \citep{cencov1982,ayJostLeSchwachhofer2015}. The square-root representation identifies the categorical Fisher metric with the metric induced from a Euclidean sphere. The distinction between affine flatness and Levi--Civita flatness is standard; global exponential or mixture coordinates do not force zero Riemannian curvature \citep{calinUdriste2014}.

Pullback geometry in learned latent spaces is also established. \citet{arvanitidisEtAl2022} pull Fisher-Rao geometry through probabilistic decoders. \citet{zhang2025} explicitly connects a language model's hidden-to-predictive map, likelihood training, and pullback geometry. For softmax representations, \citet{parkEtAl2026} develop an information-geometric account of probing and steering, while \citet{wangZhao2026} pull the output Fisher metric through transformer layers and show large deviations from ambient Euclidean geometry. \citet{sarfatiEtAl2026} empirically recover curved representation manifolds for controlled posterior beliefs and use field-aware interventions to remain on the belief family. \citet{mabrok2026} instead proposes a Fisher-metric latent semantic manifold partitioned into token regions and studies distortion induced by vocabulary discretization. These results overlap the decoder/manifold motivation and the case for geometry-aware interventions, but they do not supply the predictive-agreement stability implications proved below.

Cross-model representation convergence has been studied from several directions. The Platonic Representation Hypothesis proposes convergence toward shared statistical structure \citep{huhEtAl2024}. \citet{yuEtAl2025} align distinct neural latent geometries with pullback metrics and relative geodesic representations. \citet{huNiuVarma2026} model contextual transformations as vector fields and report shared displacement structure across language-model families. Path dependence and holonomy have also been proposed as representation diagnostics \citep{sevetlidisPavlidis2026}. The stability analysis below asks when predictive agreement forces agreement of connection-dependent transport.

Relative to all of these, the contribution here is the quantitative stability and separation package: square-root $C^2$ control of Levi--Civita transport with constants carrying no explicit vocabulary-size or token-probability factors, a paired Bernoulli construction showing that no fixed nonzero-$\alpha$ connection admits the same control without third-score-moment information, an integrated-KL-plus-Sobolev route from training risk to transport agreement, and the conditional-variance obstruction to semantic descent. To our knowledge, none of the works above states these stability or separation theorems; they supply the geometric setting and the empirical motivation rather than the agreement theory.

Continuity of Fisher information requires control of both distributions and their derivatives; \citet{rezakhaniHassaniAlipour2019} prove such bounds in classical and quantum settings. Quantitative relaxations of sufficient statistics based on bi-Lipschitz Fisher metrics are studied by \citet{yamaguchiNozawa2024}. Here a fixed comparison diffeomorphism aligns the predictive maps; their closeness is then propagated through every Amari $\alpha$-connection to an explicit parallel-transport commutator bound. Sobolev regularity supplies the derivative control needed to derive this closeness from KL convergence.

\section{Predictive statistical manifolds}

\subsection{The categorical simplex}

Let
\begin{equation}
    \simplex=\left\{p\in\R^N:p_a>0,\ \sum_{a=1}^N p_a=1\right\}.
\end{equation}
For tangent vectors $u,v,w$ satisfying $\sum_a u_a=\sum_a v_a=\sum_a w_a=0$, define
\begin{align}
    g_p^{\Delta}(u,v)&=\sum_{a=1}^N\frac{u_av_a}{p_a},\label{eq:simplex-fisher}\\
    C_p^{\Delta}(u,v,w)&=\sum_{a=1}^N\frac{u_av_aw_a}{p_a^2}.\label{eq:simplex-cubic}
\end{align}
These are the Fisher metric and Amari-Chentsov tensor. A smooth predictive map
\begin{equation}
    F:M\longrightarrow\simplex
\end{equation}
induces
\begin{equation}
    g_F=F^*g^{\Delta},\qquad C_F=F^*C^{\Delta}.
\end{equation}
When $F$ is an immersion, $g_F$ is a Riemannian metric (Lemma~\ref{lem:pullback-nondegenerate}). In local coordinates $x=(x^1,\ldots,x^m)$, writing $p_a=F_a(x)$,
\begin{align}
    (g_F)_{ij}&=\sum_a\frac{\partial_i p_a\,\partial_j p_a}{p_a},\label{eq:pullback-metric}\\
    (C_F)_{ijk}&=\sum_a\frac{\partial_i p_a\,\partial_j p_a\,\partial_k p_a}{p_a^2}.\label{eq:pullback-cubic}
\end{align}

Raise the final index of $C_F$ using $g_F$ and define $K_F$ by
\begin{equation}
    g_F(K_F(X,Y),Z)=C_F(X,Y,Z).
\end{equation}
We use the convention
\begin{equation}
    \nabla_F^{(\alpha)}=\nabla_F^{\LC}-\frac{\alpha}{2}K_F.
    \label{eq:alpha-connection}
\end{equation}
The case $\alpha=0$ is Levi-Civita by definition. That $\alpha=1$ and $\alpha=-1$ are the exponential and mixture connections on a regular exponential family is not a consequence of \eqref{eq:alpha-connection} alone; it is proved in Lemma~\ref{lem:alpha-identification}.

\subsection{Predictive quotients}

Neural predictive maps are commonly nonidentifiable. Define
\begin{equation}
    h\sim_F h'\quad\Longleftrightarrow\quad F(h)=F(h').
\end{equation}
We use a smooth predictive quotient only when it is part of the declared
structure: namely, when there are a manifold $M_F$, a surjective submersion
$\pi:M\to M_F$, and a smooth embedding $\bar F:M_F\to\simplex$ such that
$F=\bar F\circ\pi$ and the fibers of $\pi$ are exactly the equivalence classes
of $\sim_F$. Without this hypothesis we make no global quotient claim. Local
estimates may instead be applied on an explicitly chosen transverse immersed
chart on which $F$ is an immersion; an arbitrary nonintegrable horizontal
subbundle is not, by itself, a parameter manifold with a Levi--Civita
connection. Remark~\ref{rem:quotient} records this scope convention.

For a linear softmax decoder,
\begin{equation}
    F(h)=\operatorname{softmax}(Wh+b),
\end{equation}
the null space is
\begin{equation}
    \mathcal N=\{u:Wu\in\operatorname{span}\{\bm 1\}\}.
\end{equation}
The quotient $\R^d/\mathcal N$ embeds smoothly into the simplex (Lemma~\ref{lem:affine-head-structure}; an injective immersion need not be an embedding, so this requires proof), and before quotienting
\begin{equation}
    G(h)=W^\top\bigl(\operatorname{diag}p_h-p_hp_h^\top\bigr)W,
    \label{eq:softmax-metric}
    \qquad \ker G(h)=\mathcal N.
\end{equation}
The kernel identification is Lemma~\ref{lem:affine-head-structure}(i). This is the rigorous sense in which the predictively identifiable component of a hidden space is a statistical manifold. A finite set of naturally occurring hidden states remains a point cloud unless a smooth intervention chart or interpolation model is supplied.

\subsection{Ordinary addition selects a connection}

For the affine softmax family
\begin{equation}
    p_h(a)=\exp\bigl(w_a^\top h+b_a-A(h)\bigr),
\end{equation}
the quotient $\R^d/\mathcal N$ has minimal exponential-affine coordinates, so
constant quotient-coordinate vectors are $\nabla^{(1)}$-parallel by
Lemma~\ref{lem:alpha-identification}. Equivalently, ambient hidden-vector reuse
is well defined modulo $\mathcal N$. More explicitly,
\begin{equation}
    p_{h_r+h_q-h_p}(a)
    =\frac{p_{h_r}(a)p_{h_q}(a)/p_{h_p}(a)}
    {\sum_b p_{h_r}(b)p_{h_q}(b)/p_{h_p}(b)}.
    \label{eq:e-analogy}
\end{equation}
Thus hidden vector arithmetic in the ambient affine-head intervention family is exactly exponential transport, not the absence of geometry. Levi--Civita transport instead preserves Fisher lengths and angles, while mixture transport preserves the dual affine structure. Naturally occurring post-normalization states can occupy a lower-dimensional subset that is not closed under this arithmetic. Which connection best models a semantic operation on a declared intervention chart is an empirical question.

\section{Exact predictive naturality}

Exact predictive equivalence determines the corresponding metric, tensors, connections, and transports.

\begin{theorem}[Exact predictive naturality]\label{thm:exact-naturality}
Let $F_A:M_A\to\simplex$ and $F_B:M_B\to\simplex$ be smooth embeddings onto the same embedded predictive image $P$. Define
\begin{equation}
    \Phi=F_B^{-1}\circ F_A:M_A\longrightarrow M_B.
\end{equation}
Then
\begin{equation}
    \Phi^*g_B=g_A,\qquad \Phi^*C_B=C_A.
\end{equation}
For every $\alpha\in\R$,
\begin{equation}
    \dd\Phi\bigl(\nabla_A^{(\alpha)}{}_X Y\bigr)
    =\nabla_B^{(\alpha)}{}_{\dd\Phi X}(\dd\Phi Y).
    \label{eq:connection-naturality}
\end{equation}
For every piecewise $C^1$ path $\gamma$,
\begin{equation}
    \dd\Phi_{\gamma(1)}P_\gamma^{A,\alpha}
    =P_{\Phi\circ\gamma}^{B,\alpha}\dd\Phi_{\gamma(0)}.
    \label{eq:transport-naturality}
\end{equation}
Curvature tensors intertwine, and holonomies of corresponding loops are conjugate.
\end{theorem}

\begin{proof}
Since $F_A=F_B\circ\Phi$, functoriality of pullback gives
\begin{equation}
    \Phi^*g_B=\Phi^*F_B^*g^\Delta=(F_B\circ\Phi)^*g^\Delta=F_A^*g^\Delta=g_A,
\end{equation}
and identically for $C$. A Riemannian isometry preserves the Levi-Civita connection. Preservation of $g$ and $C$ implies preservation of $K$, so \eqref{eq:alpha-connection} gives \eqref{eq:connection-naturality}. Applying $\dd\Phi$ to an $A$-parallel field and using uniqueness of the linear transport ODE gives \eqref{eq:transport-naturality}. Curvature naturality follows from its definition as the commutator of covariant derivatives, and loop holonomies are consequently conjugate.
\end{proof}

\begin{corollary}[Intrinsic semantic fields]\label{cor:intrinsic-field}
Let $S_t:P\to P$ be a smooth semantic flow with infinitesimal field
\begin{equation}
    Z_P(p)=\left.\frac{\dd}{\dd t}\right|_{t=0}S_t(p).
\end{equation}
Define model fields $X_i$ by $\dd F_i(X_i)=Z_P\circ F_i$. Then
\begin{equation}
    \dd\Phi(X_A)=X_B\circ\Phi.
\end{equation}
\end{corollary}

\begin{proof}
Differentiate $F_A=F_B\circ\Phi$ and use the definitions of the two lifted fields:
\begin{equation}
\begin{split}
    \dd F_B\bigl(\dd\Phi(X_A)\bigr)
    &=\dd F_A(X_A)
      =Z_P\circ F_A\\
    &=Z_P\circ F_B\circ\Phi
      =\dd F_B(X_B\circ\Phi).
\end{split}
\end{equation}
Because $F_B$ is an embedding, $\dd F_B$ is injective at every point. Hence $\dd\Phi(X_A)=X_B\circ\Phi$.
\end{proof}

The condition is substantive: a context operation need not descend to $P$. If two hidden states have the same prediction before an intervention but distinct predictive effects afterward, then no intrinsic field $Z_P$ exists on the predictive quotient. Section~\ref{sec:semantic-obstruction} quantifies this obstruction.

\section{Prediction convergence does not imply geometric convergence}

Let $P(\cdot\mid c)$ be the data conditional and $q(\cdot\mid c)$ a model. Population cross-entropy decomposes as
\begin{equation}
    \mathcal L(q)-H_P(Y\mid C)
    =\E_C\KL\bigl(P(\cdot\mid C)\Vert q(\cdot\mid C)\bigr).
    \label{eq:ce-decomposition}
\end{equation}
The logarithmic score is strictly proper \citep{gneitingRaftery2007}. If models $q_A,q_B$ have excess risks $\epsilon_A,\epsilon_B$, then Pinsker's inequality, the triangle inequality, and Jensen's inequality for $\sqrt{\cdot}$ yield (Lemma~\ref{lem:two-model-l1})
\begin{equation}
    \E_C\|q_A(\cdot\mid C)-q_B(\cdot\mid C)\|_1
    \le \sqrt{2\epsilon_A}+\sqrt{2\epsilon_B}.
    \label{eq:two-model-pinsker}
\end{equation}
This makes output agreement natural under realizability and global population-risk optimization. It provides no derivative control.

\begin{proposition}[Regular KL-to-geometry counterexample]\label{prop:counterexample}
There exist smooth, uniformly interior, regular Bernoulli predictive families $p_n,p_*$ on $[-1,1]$ such that
\begin{equation}
    \sup_x\KL(p_*(x)\Vert p_n(x))=O(n^{-2}),
\end{equation}
    but $g_{p_n}(0)$ does not converge to $g_{p_*}(0)$.
\end{proposition}

\begin{proof}
Let $p_*(x)=\operatorname{Ber}(r(x))$ and $p_n(x)=\operatorname{Ber}(r_n(x))$, where
\begin{equation}
    r(x)=\frac13+\frac{x}{20},\qquad
    r_n(x)=r(x)+\frac{1}{100n}\sin(nx).
\end{equation}
All probabilities remain in a fixed compact subset of $(0,1)$. Moreover,
\begin{equation}
    r_n'(x)=\frac1{20}+\frac1{100}\cos(nx)\in[0.04,0.06],
\end{equation}
so every map is a regular immersion with a uniform rank bound. Uniform interiority and Taylor's theorem give $\sup_x\KL(p_*\Vert p_n)=O(n^{-2})$. The Bernoulli pullback metric is
\begin{equation}
    g(x)=\frac{r'(x)^2}{r(x)(1-r(x))}\,\dd x^2.
\end{equation}
At $x=0$, $r_n(0)=r(0)=1/3$, but $r_n'(0)=3/50$ and $r'(0)=1/20$. Hence
\begin{equation}
    g_n(0)=\frac{81}{5000}\,\dd x^2
    \neq \frac9{800}\,\dd x^2=g_*(0).
\end{equation}
\end{proof}

This example rules out an unconditional theorem of the form ``cross-entropy convergence implies Fisher convergence.'' The oscillations have bounded first derivative but unbounded higher regularity, anticipating the Sobolev condition in Section~\ref{sec:risk-to-transport}.

There is a second nonidentifiability. For any hidden diffeomorphism $\psi$,
\begin{equation}
    H'=\psi\circ H,\qquad F'=F\circ\psi^{-1}
\end{equation}
produces exactly the same predictions while changing ambient Euclidean coordinates. Pullback Fisher geometry handles this gauge covariantly, but loss alone cannot identify a privileged hidden Euclidean metric.

\section{Stability of Fisher geometry and transport}\label{sec:stability}

We now quantify how derivative-level predictive agreement propagates to geometry. We state the result on a common Euclidean chart. For two latent manifolds related by a comparison diffeomorphism $\Phi$, replace the second map by $\widetilde p=F_B\circ\Phi$.

Let $K\Subset U\subset\R^m$ and let $p,\widetilde p:U\to\simplex$. For an integer $k$, use the componentwise norm
\begin{equation}
    \|p\|_{C^k(K)}=\max_{a,|\beta|\le k}\sup_{x\in K}|\partial^\beta p_a(x)|.
\end{equation}

\begin{assumption}[Uniform regularity]\label{ass:regularity}
For constants $\eta,\lambda,B_2>0$,
\begin{align}
    p_a(x),\widetilde p_a(x)&\ge\eta,\\
    \|p\|_{C^2(K)}+\|\widetilde p\|_{C^2(K)}&\le B_2,\\
    \lambda_{\min}(g_p(x)),\lambda_{\min}(g_{\widetilde p}(x))&\ge\lambda
\end{align}
for every $a$ and $x\in K$.
\end{assumption}

The probability floor excludes the singular simplex boundary. The spectral floor excludes decoder-null directions after quotienting and prevents inverse-metric instability.

\begin{theorem}[$C^2$ stability of $\alpha$-transport]\label{thm:c2-stability}
Under Assumption~\ref{ass:regularity}, fix $A>0$. There is a constant
\begin{equation}
    C_\Gamma=C_\Gamma(m,N,A,\eta^{-1},\lambda^{-1},B_2)
\end{equation}
such that, uniformly for $|\alpha|\le A$,
\begin{equation}
\begin{split}
    &\|g_p-g_{\widetilde p}\|_{C^1(K)}
    +\|C_p-C_{\widetilde p}\|_{C^0(K)}\\
    &\hspace{3em}
    +\|\Gamma_p^{(\alpha)}-\Gamma_{\widetilde p}^{(\alpha)}\|_{C^0(K)}
    \le C_\Gamma\|p-\widetilde p\|_{C^2(K)}.
\end{split}
\label{eq:connection-stability}
\end{equation}
Let $\gamma:[0,1]\to K$ be piecewise $C^1$, with Euclidean length $L_\gamma$. If
\begin{equation}
    M_\Gamma=\max\bigl\{\|\Gamma_p^{(\alpha)}\|_{C^0},
    \|\Gamma_{\widetilde p}^{(\alpha)}\|_{C^0}\bigr\},
\end{equation}
where the norm is the induced bilinear operator norm, then
\begin{equation}
    \|P_\gamma^{p,\alpha}-P_\gamma^{\widetilde p,\alpha}\|_{\op}
    \le L_\gamma e^{M_\Gamma L_\gamma}C_\Gamma
    \|p-\widetilde p\|_{C^2(K)}.
    \label{eq:pt-stability}
\end{equation}
\end{theorem}

\begin{proof}
Equations \eqref{eq:pullback-metric} and \eqref{eq:pullback-cubic} are finite sums of rational functions of $p$ and its first derivatives. Their denominators are controlled by $\eta$. Differentiating \eqref{eq:pullback-metric} once uses at most second derivatives, so the first two terms of \eqref{eq:connection-stability} follow from the mean-value inequality on the compact \emph{convex} regularity class; Lemma~\ref{lem:lipschitz-class} carries this out and records why convexity is needed.

Lemma~\ref{lem:lower-alpha} gives the lower-index $\alpha$-coefficients directly
in probability coordinates:
\begin{equation}
    \Gamma^{(\alpha)}_{ij,k}
    =\sum_a\left[
    \frac{\partial_{ij}p_a\,\partial_kp_a}{p_a}
    -\frac{1+\alpha}{2}
    \frac{\partial_ip_a\,\partial_jp_a\,\partial_kp_a}{p_a^2}
    \right].
    \label{eq:lower-alpha}
\end{equation}
Thus their difference is $O(\|p-\widetilde p\|_{C^2})$. Raising the final index requires inverse metrics. The identity
\begin{equation}
    g_p^{-1}-g_{\widetilde p}^{-1}
    =g_p^{-1}(g_{\widetilde p}-g_p)g_{\widetilde p}^{-1}
\end{equation}
and the spectral floor $\lambda$ give inverse-metric stability, proving \eqref{eq:connection-stability}.

Along $\gamma$, transport solves the linear ODE
\begin{equation}
    \dot V(t)=-\Gamma^{(\alpha)}(\gamma(t))[\dot\gamma(t)]V(t).
\end{equation}
Duhamel's formula bounds the difference of the two fundamental solutions by the integral of the coefficient defect multiplied by their growth factors. Gronwall's inequality then yields \eqref{eq:pt-stability}; the two growth factors multiply to a constant along the path, which is what produces the single exponential rather than $e^{2M_\Gamma L_\gamma}$ (Lemma~\ref{lem:duhamel-cancellation}).
\end{proof}

\begin{corollary}[Cross-model transport commutator]\label{cor:cross-model-commutator}
Let $F_A:M_A\to\simplex$, $F_B:M_B\to\simplex$, and let $\Phi:M_A\to M_B$ be a $C^2$ local diffeomorphism. Suppose $F_A$ and $F_B\circ\Phi$ satisfy Assumption~\ref{ass:regularity} and
\begin{equation}
    \|F_A-F_B\circ\Phi\|_{C^2(K)}\le\varepsilon.
\end{equation}
Then
\begin{equation}
\begin{split}
    &\bigl\|\dd\Phi_{\gamma(1)}P_\gamma^{A,\alpha}
    -P_{\Phi\circ\gamma}^{B,\alpha}\dd\Phi_{\gamma(0)}\bigr\|_{\op}\\
    &\hspace{5em}\le L_\gamma e^{M_\Gamma L_\gamma}C_\Gamma\varepsilon,
\end{split}
\label{eq:commutator-bound}
\end{equation}
where the operator norm measures inputs with the background norm on $T_{\gamma(0)}M_A$ and outputs with the \emph{aligned} tangent norm on $T_{\Phi(\gamma(1))}M_B$, defined as the pullback of the background norm through the alignment:
\begin{equation}
    \|w\|_{\mathrm{al},x}:=\|\dd\Phi_x^{-1}w\|,
    \qquad w\in T_{\Phi(x)}M_B .
\end{equation}
In these norms the left-hand side equals
$\|P_\gamma^{A,\alpha}-\dd\Phi_{\gamma(1)}^{-1}P_{\Phi\circ\gamma}^{B,\alpha}\dd\Phi_{\gamma(0)}\|_{\op}$
on $M_A$, which is the quantity the proof bounds. If an independently chosen background norm is used on $TM_B$ instead, the right-hand side acquires the additional factor $\sup_{x\in K}\|\dd\Phi_x\|_{\op}$, which must then be stated. At $\varepsilon=0$, this reduces to Theorem~\ref{thm:exact-naturality}.
\end{corollary}

\begin{proof}
The connection induced by $F_B\circ\Phi$ is the pullback of the connection induced by $F_B$. Its transport is conjugated by $\dd\Phi$. Apply Theorem~\ref{thm:c2-stability} on $M_A$.
\end{proof}

\begin{theorem}[$C^3$ curvature stability]\label{thm:c3-curvature}
Assume the hypotheses of Theorem~\ref{thm:c2-stability}, replace the $C^2$ bound by
\begin{equation}
    \|p\|_{C^3(K)}+\|\widetilde p\|_{C^3(K)}\le B_3,
\end{equation}
and retain the probability and spectral floors. Then
\begin{equation}
    \|R_p^{(\alpha)}-R_{\widetilde p}^{(\alpha)}\|_{C^0(K)}
    \le C_R\|p-\widetilde p\|_{C^3(K)},
\end{equation}
where $C_R$ depends on $m,N,A,\eta^{-1},\lambda^{-1},B_3$. Levi-Civita sectional curvatures also converge on corresponding two-planes whose Gram determinants are uniformly bounded below.
\end{theorem}

\begin{proof}
The lower-index Christoffel coefficients are smooth rational functions of the
second jet $J^2p=(p,Dp,D^2p)$ whose only denominators are powers of $p_a$.
Their first derivatives are smooth rational functions of $J^3p$. The
mean-value inequality on the convex probability-and-derivative envelope, where
$p_a\ge\eta$, therefore controls their $C^1$ difference by
$C\|p-\widetilde p\|_{C^3}$. Raising an index does not require a line segment
through the nonconvex spectral-floor class: use directly
$g^{-1}-\widetilde g^{-1}=g^{-1}(\widetilde g-g)\widetilde g^{-1}$ and
$\partial g^{-1}=-g^{-1}(\partial g)g^{-1}$, expand the difference of the
latter identity, and apply the endpoint bounds
$\|g^{-1}\|,\|\widetilde g^{-1}\|\le\lambda^{-1}$. This gives
\begin{equation}
  \|\Gamma_p^{(\alpha)}-\Gamma_{\widetilde p}^{(\alpha)}\|_{C^1(K)}
  \le C\|p-\widetilde p\|_{C^3(K)}.
\end{equation}
The same compact jet class supplies uniform $C^1$ bounds for each connection.
Insert these estimates into
\begin{equation}
    R^\ell{}_{ijk}=\partial_i\Gamma^\ell{}_{jk}-\partial_j\Gamma^\ell{}_{ik}
    +\Gamma^\ell{}_{ir}\Gamma^r{}_{jk}-\Gamma^\ell{}_{jr}\Gamma^r{}_{ik}.
\end{equation}
The displayed expression is polynomial in $\partial\Gamma$ and quadratic in
$\Gamma$, so subtraction and the product inequality prove the asserted
curvature estimate. For sectional curvature, choose background-orthonormal
representatives $u,v$ of each compared plane. The numerator
$g(R(u,v)v,u)$ is controlled by the curvature and metric estimates, while the
assumed positive lower bound for
$g(u,u)g(v,v)-g(u,v)^2$ controls the denominator. This proves the final claim.
\end{proof}

\subsection{Boundary-robust Levi--Civita stability}

The preceding estimates are uniform over a probability-coordinate regularity class and are intentionally general in $\alpha$. For Levi--Civita geometry, the square-root representation gives a sharper result. Define
\begin{equation}
    \psi_p(x)=2\bigl(\sqrt{p_1(x)},\ldots,\sqrt{p_N(x)}\bigr)
    \in\R^N.
    \label{eq:square-root-map}
\end{equation}
Its image lies on the sphere of radius two. Equip target-valued derivatives with their Euclidean operator norms, and put
\begin{equation}
    \|D^j\psi\|_{C^0(K),\op}
    =\sup_{\substack{x\in K\\|u_1|,\ldots,|u_j|\le1}}
      \|D^j\psi_x[u_1,\ldots,u_j]\|_{\ell_2},
\end{equation}
and put
\begin{equation}
    \delta_j=\sup_{x\in K}\|D^j\psi_p(x)-D^j\psi_{\widetilde p}(x)\|_{\op},
    \qquad j=1,2.
\end{equation}

\begin{theorem}[Vocabulary-independent Levi--Civita stability]\label{thm:sqrt-lc-stability}
Suppose $\psi_p,\psi_{\widetilde p}:U\to\R^N$ are $C^2$ immersions and, on $K$,
\begin{align}
    \|D\psi_p\|_{\op},\|D\psi_{\widetilde p}\|_{\op}&\le B_1,\nonumber\\
    \|D^2\psi_p\|_{\op},\|D^2\psi_{\widetilde p}\|_{\op}&\le B_2,\label{eq:sqrt-regularity}\\
    g_p,g_{\widetilde p}&\succeq\lambda I.\nonumber
\end{align}
Then
\begin{align}
    \|g_p-g_{\widetilde p}\|_{C^0(K)}
    &\le 2B_1\delta_1,\label{eq:sqrt-metric-bound}\\
    \|Dg_p-Dg_{\widetilde p}\|_{C^0(K)}
    &\le2(B_1\delta_2+B_2\delta_1),\label{eq:sqrt-metric-derivative-bound}\\
    \|\Gamma_p^{\LC}-\Gamma_{\widetilde p}^{\LC}\|_{C^0(K)}
    &\le D_{\LC},\label{eq:sqrt-connection-bound}
\end{align}
where
\begin{equation}
    D_{\LC}
    =\lambda^{-1}(B_1\delta_2+B_2\delta_1)
    +2B_1^2B_2\lambda^{-2}\delta_1.
    \label{eq:dlc}
\end{equation}
For a path $\gamma$ of Euclidean length $L_\gamma$,
\begin{equation}
    \|P_\gamma^{p,\LC}-P_\gamma^{\widetilde p,\LC}\|_{\op}
    \le L_\gamma D_{\LC}
    \min\!\left\{
        \exp\!\left(\frac{B_1B_2}{\lambda}L_\gamma\right),
        \frac{B_1^2}{\lambda}
    \right\}.
    \label{eq:sqrt-transport-bound}
\end{equation}
Every displayed constant is independent of $N$ and of $\min_a p_a,\min_a\widetilde p_a$. The second transport branch uses metric compatibility and prevents an unavoidable exponential-in-length interpretation. It does not remove rank conditioning: after substituting $D_{\LC}$, that branch contains terms proportional to $\lambda^{-2}$ and $\lambda^{-3}$.
\end{theorem}

\begin{proof}
Differentiating \eqref{eq:square-root-map} gives the exact identities
\begin{equation}
    (g_p)_{ij}=\langle\partial_i\psi_p,\partial_j\psi_p\rangle,
    \qquad
    (\Gamma_p^{\LC})_{ij,k}
    =\langle\partial_{ij}\psi_p,\partial_k\psi_p\rangle.
    \label{eq:induced-lc-identities}
\end{equation}
The first is \eqref{eq:pullback-metric}; the second follows either from the Koszul formula or by differentiating the first and cancelling the symmetric terms. Subtracting the two Gram matrices proves \eqref{eq:sqrt-metric-bound}. Differentiating them once and applying the product rule proves \eqref{eq:sqrt-metric-derivative-bound}. The lower-index coefficients in \eqref{eq:induced-lc-identities} differ by at most $B_1\delta_2+B_2\delta_1$. Moreover,
\begin{equation}
    \|g_p^{-1}-g_{\widetilde p}^{-1}\|_{\op}
    \le 2B_1\lambda^{-2}\delta_1,
\end{equation}
and each lower-index coefficient has norm at most $B_1B_2$. Raising the final index yields \eqref{eq:sqrt-connection-bound}--\eqref{eq:dlc}. The same identities give $\|\Gamma^{\LC}\|_{\op}\le B_1B_2/\lambda$, so Duhamel's formula and Gronwall's inequality give the exponential branch of \eqref{eq:sqrt-transport-bound}.

For the second branch, note that
\begin{equation}
    \lambda I\preceq g_p,g_{\widetilde p}\preceq B_1^2I.
\end{equation}
Levi--Civita transport is an isometry for its own metric; hence every subpath propagator, measured in the background Euclidean norm, has operator norm at most $B_1/\sqrt\lambda$. The exact variation-of-connection identity contains one propagator from each geometry and one copy of their connection defect. Taking norms and integrating $|\dot\gamma|$ gives
\begin{equation}
    \|P_\gamma^{p,\LC}-P_\gamma^{\widetilde p,\LC}\|_{\op}
    \le L_\gamma\frac{B_1^2}{\lambda}D_{\LC}.
\end{equation}
Taking the better of the two valid estimates proves \eqref{eq:sqrt-transport-bound}. Appendix~\ref{app:sqrt-audit} gives the full Duhamel identity and its finite-scale interpretation.
\end{proof}

\begin{corollary}[Boundary-robust Levi--Civita curvature]\label{cor:sqrt-lc-curvature}
Under the hypotheses of Theorem~\ref{thm:sqrt-lc-stability}, suppose additionally that the two square-root immersions are $C^3$ so that their classical curvature tensors are defined. Then the covariant Levi--Civita curvature tensors satisfy
\begin{equation}
    \|\operatorname{Rm}_p-\operatorname{Rm}_{\widetilde p}\|_{C^0(K)}
    \le C(B_1,B_2,\lambda^{-1})(\delta_1+\delta_2).
\end{equation}
Consequently, sectional curvatures converge on corresponding planes whose Fisher Gram determinants are uniformly bounded below. The constant is again independent of vocabulary size and minimum token probability.
\end{corollary}

\begin{proof}
Regard each square-root map as a Euclidean immersion. Its second fundamental form is
\begin{equation}
    \mathrm{II}_p(u,v)
    =(I-\Pi_p)D^2\psi_p(u,v),
    \qquad
    \Pi_p=D\psi_p\,g_p^{-1}(D\psi_p)^*.
\end{equation}
The inverse-metric estimate used above controls $\Pi_p-\Pi_{\widetilde p}$ by $C(B_1,\lambda^{-1})\delta_1$, and hence controls the second fundamental forms by $C(B_1,B_2,\lambda^{-1})(\delta_1+\delta_2)$. The Euclidean Gauss equation expresses $\operatorname{Rm}$ as a quadratic contraction of $\mathrm{II}$, proving the tensor bound. The sectional-curvature statement follows by controlling its metric denominator.
\end{proof}

\subsection{Score moments and the affine-head specialization}

The absence of explicit boundary factors above is not an algebraic accident. Let $Y\sim p(x)$ and define the score and its derivative by
\begin{equation}
    S_i(Y,x)=\partial_i\log p_Y(x),
    \qquad H_{ij}(Y,x)=\partial_iS_j(Y,x).
\end{equation}

\begin{proposition}[Score-moment form of the $\alpha$-connection]\label{prop:score-moments}
For every smooth categorical family with nondegenerate pullback Fisher metric,
\begin{align}
    (g_p)_{ij}&=\E_p[S_iS_j],\label{eq:score-metric}\\
    (C_p)_{ijk}&=\E_p[S_iS_jS_k],\label{eq:score-cubic}\\
    \Gamma^{(\alpha)}_{ij,k}
    &=\E_p[H_{ij}S_k]
      +\frac{1-\alpha}{2}\E_p[S_iS_jS_k].
      \label{eq:score-alpha}
\end{align}
For two affine softmax heads expressed on a common \emph{minimal}
natural-coordinate chart---equivalently, after quotienting their predictive
null spaces---let $G,C$ and $\widetilde G,\widetilde C$ denote the corresponding
second and third centered moments of the decoder rows. Then $G$ and
$\widetilde G$ are positive definite and
\begin{equation}
    \Gamma^{(\alpha)}=\frac{1-\alpha}{2}G^{-1}C.
    \label{eq:affine-alpha-moment}
\end{equation}
If $G,\widetilde G\succeq\lambda I$, then
\begin{equation}
\begin{split}
    \|\Gamma^{(\alpha)}-\widetilde\Gamma^{(\alpha)}\|_{\op}
    \le\frac{|1-\alpha|}{2}\Bigl(&\lambda^{-1}\|C-\widetilde C\|_{\op}\\
    &+\lambda^{-2}\|G-\widetilde G\|_{\op}
      \|\widetilde C\|_{\op}\Bigr).
\end{split}
\label{eq:affine-alpha-stability}
\end{equation}
This bound has no explicit dependence on vocabulary size or minimum token probability.
\end{proposition}

\begin{proof}
Since $\partial_i p_a=p_aS_i(a)$ and
\begin{equation}
    \partial_{ij}p_a=p_a\bigl(H_{ij}(a)+S_i(a)S_j(a)\bigr),
\end{equation}
substitution into \eqref{eq:pullback-metric}, \eqref{eq:pullback-cubic}, and
\eqref{eq:lower-alpha} (Lemma~\ref{lem:lower-alpha}) proves
\eqref{eq:score-metric}--\eqref{eq:score-alpha}. At an affine head,
$S(Y)=w_Y-\E_pw_Y$ and $H_{ij}(Y)=-G_{ij}$ is independent of $Y$.
Because $\E_pS=0$, the first term of \eqref{eq:score-alpha} vanishes. On the
declared minimal chart $G$ is invertible, so raising the final index gives
\eqref{eq:affine-alpha-moment}. Finally use
\begin{equation}
    G^{-1}C-\widetilde G^{-1}\widetilde C
    =G^{-1}(C-\widetilde C)
     +(G^{-1}-\widetilde G^{-1})\widetilde C
\end{equation}
and $\|G^{-1}-\widetilde G^{-1}\|_{\op}
\le\lambda^{-2}\|G-\widetilde G\|_{\op}$.
\end{proof}

\begin{corollary}[General $\alpha$ stability from cubic-moment control]\label{cor:alpha-moment-stability}
In the setting of Theorem~\ref{thm:sqrt-lc-stability}, suppose additionally that
\begin{equation}
    \|C_p\|_{C^0(K),\op},\|C_{\widetilde p}\|_{C^0(K),\op}\le M_3,
    \qquad
    \|C_p-C_{\widetilde p}\|_{C^0(K),\op}\le\delta_C.
\end{equation}
Then
\begin{equation}
\begin{split}
    \|\Gamma_p^{(\alpha)}-\Gamma_{\widetilde p}^{(\alpha)}\|_{\op}
    \le D_{\LC}+\frac{|\alpha|}{2}\left(
        \lambda^{-1}\delta_C
        +2B_1M_3\lambda^{-2}\delta_1
    \right).
\end{split}
\label{eq:general-alpha-moment-bound}
\end{equation}
Thus square-root $C^2$ control alone determines Levi--Civita stability; a general nonzero-$\alpha$ connection additionally requires stability of the raised third score moment. A tokenwise probability floor is one sufficient route, not a necessary one. Affine softmax structure supplies a floor-free route directly through Proposition~\ref{prop:score-moments}.
\end{corollary}

\begin{proof}
Use $\Gamma^{(\alpha)}=\Gamma^{\LC}-(\alpha/2)g^{-1}C$, Theorem~\ref{thm:sqrt-lc-stability}, and the inverse-metric estimate in its proof.
\end{proof}

\begin{proposition}[Square-root control does not control nonzero $\alpha$]\label{prop:boundary-alpha-separation}
There is a one-dimensional categorical family whose square-root map has uniformly bounded derivatives of every order and whose Fisher metric is identically one, while every fixed nonzero-$\alpha$ connection coefficient diverges toward the simplex boundary.
\end{proposition}

\begin{proof}
For $0<\theta<\pi$, take the Bernoulli family
\begin{equation}
    p_\theta
    =\left(\sin^2\frac{\theta}{2},
            \cos^2\frac{\theta}{2}\right).
\end{equation}
Its square-root representation is
\begin{equation}
    \psi_\theta
    =2\left(\sin\frac{\theta}{2},
             \cos\frac{\theta}{2}\right).
\end{equation}
All derivatives are bounded as $\theta\downarrow0$, while
\begin{equation}
    g_{\theta\theta}=1,
    \qquad
    \Gamma^{\LC}_{\theta\theta,\theta}
    =\langle\psi_\theta'',\psi_\theta'\rangle=0.
\end{equation}
The two score values are $\cot(\theta/2)$ and $-\tan(\theta/2)$, so
\begin{equation}
    C_{\theta\theta\theta}=2\cot\theta.
\end{equation}
Consequently,
\begin{equation}
    \Gamma^{(\alpha)}_{\theta\theta,\theta}
    =-\alpha\cot\theta,
\end{equation}
which diverges for every $\alpha\ne0$. Thus square-root $C^2$ control and a perfect Fisher spectral floor are sufficient for Levi--Civita stability but not for a general $\alpha$-connection. At an affine softmax head in natural coordinates, Proposition~\ref{prop:score-moments} supplies a separate moment-based route and $\Gamma^{(1)}=0$ identically.
\end{proof}

\begin{corollary}[Sharp connection selectivity of square-root control]\label{cor:alpha-zero-selectivity}
Within the Amari $\alpha$-family, $\alpha=0$ is the unique fixed parameter for which uniform square-root $C^2$ bounds, square-root $C^2$ agreement, and a Fisher spectral floor alone imply uniform connection stability in a declared background chart. More precisely, for every fixed $\alpha\ne0$ there are pairs of Bernoulli predictive maps on the same compact chart such that their square-root maps converge to one another in $C^2$, both Fisher metrics are identically one, and yet
\begin{equation}
    \|\Gamma_{p_n}^{(\alpha)}-
      \Gamma_{\widetilde p_n}^{(\alpha)}\|_{C^0}
    \longrightarrow\infty.
    \label{eq:nonzero-alpha-pair-divergence}
\end{equation}
This is a uniqueness statement only for stability from these stated data; it is not a claim that Levi--Civita is the unique geometrically natural connection.
\end{corollary}

\begin{proof}
The $\alpha=0$ implication is Theorem~\ref{thm:sqrt-lc-stability}. For the converse, let $K=[0,1]$, $\varepsilon_n=1/n$, and use the Bernoulli construction in Proposition~\ref{prop:boundary-alpha-separation} with angles
\begin{equation}
    \theta_n(x)=x+\varepsilon_n,
    \qquad
    \widetilde\theta_n(x)=x+2\varepsilon_n.
\end{equation}
The associated square-root maps are phase shifts of the same radius-two circle and therefore converge to one another in $C^2(K)$ at rate $O(\varepsilon_n)$; their first and second derivative envelopes are respectively one and one half, and both induced metrics are $\dd x^2$. At $x=0$ their connection-coefficient defect is
\begin{equation}
    |\alpha|\,|\cot\varepsilon_n-\cot(2\varepsilon_n)|
    \sim \frac{|\alpha|}{2\varepsilon_n},
\end{equation}
which proves \eqref{eq:nonzero-alpha-pair-divergence}.
\end{proof}

\begin{remark}[Why the assumptions matter]
The probability-coordinate constants in Theorems~\ref{thm:c2-stability} and \ref{thm:c3-curvature} diverge near $p_a=0$ because their uniform class controls derivatives of $p_a$ independently of $p_a$. This is not always an intrinsic blow-up of the pulled-back geometry: Theorem~\ref{thm:sqrt-lc-stability} captures cancellations through square-root derivatives, while Proposition~\ref{prop:score-moments} captures them through probability-weighted score moments. Rank loss remains intrinsic because raising indices requires $g^{-1}$. For nonzero $\alpha$, the raised third score moment supplies information not determined by the Fisher metric alone and must be controlled explicitly; uniform interiority is sufficient but not logically necessary. Vocabulary independence means that no explicit $N$ occurs in the displayed constants; the derivative bounds and spectral floor must still be uniform and may deteriorate with model scale. The estimates are conditional finite-scale bounds, not numerical certificates until those quantities are measured.
\end{remark}

\section{From cross-entropy to transport under Sobolev regularity}\label{sec:risk-to-transport}

Theorem~\ref{thm:c2-stability} assumes derivative-level agreement. We now give sufficient conditions under which population KL convergence implies it.

Let $\Omega$ be either a compact $m$-dimensional smooth manifold without
boundary or a bounded Sobolev-extension domain in $\R^m$. On a compact
manifold fix a smooth background Riemannian metric $h$ and define the scalar
$H^t$ scale spectrally by $(I+\Delta_h)^{t/2}$, where $\Delta_h$ is the
nonnegative Laplace--Beltrami operator; the Hilbert-valued scale is its
Hilbert tensor product with the target. On a bounded domain, ``Sobolev-extension
domain through order $s$'' means that the restriction-space Bessel-potential
norms are used and that there is one scalar linear operator
\begin{equation}
  E_\Omega:H^t(\Omega)\longrightarrow H^t(\R^m),
  \qquad 0\le t\le s,
  \label{eq:total-extension-definition}
\end{equation}
which restricts to the identity on $\Omega$ and is bounded at every indicated
order. This common-operator property is part of the hypothesis, not inferred
from the existence of unrelated extension operators at separate orders. All
volume measures, covariant $C^k$ norms, operator norms, and path lengths in this
section refer to $h$ on the manifold and to the Euclidean background on the
domain; write $h$ for the selected background metric in either case. Let $\mu$
have density $\rho$ with respect to volume, satisfying
$\rho\ge\rho_0>0$. Define
\begin{equation}
    \mathcal K(p,\widetilde p)
    =\int_\Omega \KL\bigl(p(x)\Vert\widetilde p(x)\bigr)\,\dd\mu(x).
\end{equation}

\begin{theorem}[Risk-to-transport convergence]\label{thm:risk-to-transport}
Let $p,\widetilde p:\Omega\to\simplex$ satisfy, for some $\eta,\lambda>0$,
\begin{equation}
  p_a,\widetilde p_a\ge\eta,\qquad
  g_p,g_{\widetilde p}\succeq\lambda h.
\end{equation}
Suppose
\begin{equation}
    p,\widetilde p\in H^s(\Omega;\R^N),\qquad
    s>2+\frac m2,
\end{equation}
and
\begin{equation}
    \|p\|_{H^s}+\|\widetilde p\|_{H^s}\le B.
\end{equation}
Fix $A,L_0>0$. For every $r$ with
\begin{equation}
    2+\frac m2<r<s,
\end{equation}
there is a constant $C=C(\Omega,m,s,r,N,A,L_0,\rho_0^{-1},\eta^{-1},\lambda^{-1},B)$ such that, uniformly for $|\alpha|\le A$ and every path $\gamma$ of background length $L_\gamma\le L_0$,
\begin{equation}
    \|P_\gamma^{p,\alpha}-P_\gamma^{\widetilde p,\alpha}\|_{\op}
    \le C B^{r/s}\mathcal K(p,\widetilde p)^{\frac{s-r}{2s}}.
    \label{eq:risk-to-transport}
\end{equation}
Equivalently, the transport error is $O(\mathcal K^\beta)$ for every
\begin{equation}
    0<\beta<\frac{s-2-m/2}{2s}.
\end{equation}
\end{theorem}

\begin{proof}
Pinsker's inequality and $\rho\ge\rho_0$ give
\begin{equation}
\begin{split}
    \|p-\widetilde p\|_{L^2(\dd\mathrm{vol})}^2
    &\le\rho_0^{-1}\int_\Omega\|p-\widetilde p\|_1^2\,\dd\mu\\
    &\le 2\rho_0^{-1}\mathcal K(p,\widetilde p).
\end{split}
\end{equation}
Sobolev interpolation \citep{adamsFournier2003} yields
\begin{equation}
    \|p-\widetilde p\|_{H^r}
    \le C\|p-\widetilde p\|_{L^2}^{1-r/s}
    \|p-\widetilde p\|_{H^s}^{r/s},
\end{equation}
and therefore
\begin{equation}
    \|p-\widetilde p\|_{H^r}
    \le C B^{r/s}\mathcal K(p,\widetilde p)^{(s-r)/(2s)}.
\end{equation}
Since $r>2+m/2$, Lemma~\ref{lem:hilbert-domain} gives
$H^r\hookrightarrow C^2$ in the declared background norm. On a Euclidean
extension domain, the pointwise proof of Theorem~\ref{thm:c2-stability} applies
on the compact closure. On a compact manifold, apply the covariant form,
Lemma~\ref{lem:chart-transfer}. In either case the $H^s$ envelope supplies the
required $C^2$ envelope, the probability and spectral floors bound the
individual connection coefficients, and $L_\gamma\le L_0$ absorbs the
Duhamel growth factor into the displayed constant.
\end{proof}

\begin{theorem}[Boundary-robust risk-to-Levi--Civita transport]\label{thm:sqrt-risk-to-transport}
Let $p,\widetilde p:\Omega\to\simplex$, write $\psi_p=2\sqrt p$ and $\psi_{\widetilde p}=2\sqrt{\widetilde p}$ componentwise, and assume
\begin{equation}
    \psi_p,\psi_{\widetilde p}\in H^s(\Omega;\ell_2^N),
    \qquad s>2+\frac m2,
\end{equation}
with
\begin{equation}
    \|\psi_p\|_{H^s}+\|\psi_{\widetilde p}\|_{H^s}\le B,
    \qquad g_p,g_{\widetilde p}\succeq\lambda h.
\end{equation}
Fix $L_0>0$. For every $r$ satisfying $2+m/2<r<s$ and every path of background length at most $L_0$,
\begin{equation}
    \|P_\gamma^{p,\LC}-P_\gamma^{\widetilde p,\LC}\|_{\op}
    \le C\,\mathcal K(p,\widetilde p)^{\frac{s-r}{2s}},
    \label{eq:sqrt-risk-transport}
\end{equation}
where $C$ depends on $\Omega,s,r,\rho_0^{-1},B,\lambda^{-1},L_0$ but not on $N$ or a minimum token probability. Equivalently, the transport error is $O(\mathcal K^\beta)$ for every
\begin{equation}
    0<\beta<\frac{s-2-m/2}{2s}.
\end{equation}
\end{theorem}

\begin{proof}
For categorical distributions Lemma~\ref{lem:hellinger-kl} gives the scalar inequality
\begin{equation}
    \sum_a(\sqrt{p_a}-\sqrt{\widetilde p_a})^2
    \le\KL(p\Vert\widetilde p)
\end{equation}
and hence
\begin{equation}
    \|\psi_p-\psi_{\widetilde p}\|_{L^2(\dd\mathrm{vol};\ell_2)}^2
    \le4\rho_0^{-1}\mathcal K(p,\widetilde p).
    \label{eq:sqrt-kl}
\end{equation}
Hilbert-valued Sobolev interpolation and embedding, with constants independent of the target dimension (Lemmas~\ref{lem:hilbert-interpolation}--\ref{lem:hilbert-domain}), give
\begin{equation}
\begin{split}
    \|\psi_p-\psi_{\widetilde p}\|_{C^2(\ell_2)}
    &\le C\|\psi_p-\psi_{\widetilde p}\|_{H^r(\ell_2)}\\
    &\le C B^{r/s}\mathcal K(p,\widetilde p)^{(s-r)/(2s)}.
\end{split}
\end{equation}
The same embedding bounds the first two derivatives of each square-root map in
terms of $B$. Apply Theorem~\ref{thm:sqrt-lc-stability} on a Euclidean extension
domain and Lemma~\ref{lem:manifold-sqrt-stability} on a compact manifold.
Remark~\ref{rem:where-N-free} tracks where each constant comes from and
confirms that no step introduces a factor of $N$ or of $\min_ap_a$.
\end{proof}

\begin{corollary}[Two-model geometric convergence]\label{cor:two-model-convergence}
Suppose two model sequences $p_{A,n},p_{B,n}$ on a shared intervention chart converge in population cross-entropy to the same conditional map $p_*$, and all three maps satisfy uniform versions of the hypotheses of Theorem~\ref{thm:risk-to-transport}. Then, for every admissible $\beta$,
\begin{equation}
    \|P_\gamma^{A,n,\alpha}-P_\gamma^{B,n,\alpha}\|_{\op}
    \le C\bigl(\epsilon_{A,n}^{\beta}+\epsilon_{B,n}^{\beta}\bigr),
\end{equation}
where $\epsilon_{i,n}=\E\KL(p_*\Vert p_{i,n})$. If the models instead live on aligned charts, the same statement holds for the transport commutator in \eqref{eq:commutator-bound} after replacing $p_{B,n}$ by $p_{B,n}\circ\Phi_n$, provided these pulled-back maps converge to the same $p_*$ and satisfy the same uniform hypotheses.

Alternatively, if the square-root maps satisfy uniform versions of Theorem~\ref{thm:sqrt-risk-to-transport}, their Levi--Civita transports obey the same conclusion with constants independent of vocabulary size and minimum token probability.
\end{corollary}

\begin{proof}
Let $P_\gamma^{*,\alpha}$ denote the $\alpha$-transport induced by $p_*$. For $i\in\{A,B\}$, Theorem~\ref{thm:risk-to-transport} and the uniform hypotheses give
\begin{equation}
    \|P_\gamma^{i,n,\alpha}-P_\gamma^{*,\alpha}\|_{\op}
    \le C_i\epsilon_{i,n}^{\beta}.
\end{equation}
The constants are uniformly bounded in $n$ and may be replaced by one constant $C$. The operator-norm triangle inequality, applied to the common target transport rather than to the KL divergences, now gives
\begin{equation}
\begin{split}
    \|P_\gamma^{A,n,\alpha}-P_\gamma^{B,n,\alpha}\|_{\op}
    &\le
    \|P_\gamma^{A,n,\alpha}-P_\gamma^{*,\alpha}\|_{\op}
    +\|P_\gamma^{*,\alpha}-P_\gamma^{B,n,\alpha}\|_{\op}\\
    &\le C\bigl(\epsilon_{A,n}^{\beta}+\epsilon_{B,n}^{\beta}\bigr).
\end{split}
\end{equation}
For an alignment $\Phi_n$, work on model A's chart with the pulled-back predictive map $p_{B,n}\circ\Phi_n$. Exact naturality identifies its transport with
\begin{equation}
    (\dd\Phi_n)_{\gamma(1)}^{-1}
    P_{\Phi_n\circ\gamma}^{B,n,\alpha}
    (\dd\Phi_n)_{\gamma(0)}.
\end{equation}
The same triangle argument therefore gives the commutator estimate in the aligned norms of Corollary~\ref{cor:cross-model-commutator}; an independently chosen target background norm incurs the additional alignment factor stated there. Replacing Theorem~\ref{thm:risk-to-transport} by Theorem~\ref{thm:sqrt-risk-to-transport} proves the Levi--Civita alternative.
\end{proof}

If $s>3+m/2$, the same argument combined with
Theorem~\ref{thm:c3-curvature} on a Euclidean domain and
Lemma~\ref{lem:chart-transfer} on a compact manifold gives curvature convergence
for every
\begin{equation}
    0<\beta<\frac{s-3-m/2}{2s}.
\end{equation}

For Levi--Civita curvature, Corollary~\ref{cor:sqrt-lc-curvature}, and its
manifold form in Lemma~\ref{lem:manifold-sqrt-stability}, instead use $C^2$
convergence of the square-root immersions. Thus
Theorem~\ref{thm:sqrt-risk-to-transport} also gives vocabulary-independent LC
curvature convergence at the transport exponent, provided additionally that
the square-root immersions are $C^3$ so their classical curvature tensors are
defined and that the compared plane Gram determinants stay uniformly positive.

The uniform Sobolev hypothesis excludes the high-frequency behavior in Proposition~\ref{prop:counterexample}. For language models it should be tested on low-dimensional continuous intervention charts, such as soft-prompt paths or controlled activation subspaces, rather than assumed on the discrete set of tokenized contexts.

\section{The semantic-sufficiency obstruction}\label{sec:semantic-obstruction}

Even perfect agreement of predictive geometry does not guarantee that a semantic operation is a field on that geometry. We now separate this representational obstruction from connection mismatch.

Let $F:\mathcal H\to P$ be a declared smooth predictive quotient onto a
finite-dimensional embedded submanifold $P\subset\simplex$, let $H$ be a random
hidden state with law $\mu$, and let $X_T(H)\in T_H\mathcal H$ be measurable and
describe a semantic operation $T$. Its observable infinitesimal predictive
effect is
\begin{equation}
    Z_T(H)=\dd F_H X_T(H)\in T_{F(H)}P.
\end{equation}
Assume square integrability in the Fisher metric. Since $P$ is embedded in the simplex, tangent-valued conditional expectations may be taken fiberwise in the ambient zero-sum space; conditioning on $F(H)=p$ fixes both $T_pP$ and $g_p$.

\begin{theorem}[Semantic conditional-variance obstruction]\label{thm:semantic-obstruction}
For a measurable, square-integrable field $Y:P\to TP$, define
\begin{equation}
    \mathcal E_T(Y)=\E\|Z_T(H)-Y(F(H))\|_{g_{F(H)}}^2.
\end{equation}
Let $\nu=\mathcal L(F(H))$. The conditional mean below denotes any measurable
version and is defined $\nu$-almost everywhere:
\begin{equation}
    \overline Z_T(p)=\E[Z_T(H)\mid F(H)=p].
\end{equation}
Then
\begin{equation}
\begin{split}
    \mathcal E_T(Y)
    ={}&\E\|Z_T-\overline Z_T(F(H))\|_g^2\\
    &+\E\|\overline Z_T(F(H))-Y(F(H))\|_g^2.
\end{split}
\label{eq:semantic-pythagoras}
\end{equation}
Consequently,
\begin{equation}
    \inf_Y\mathcal E_T(Y)
    =\E\Var_g(Z_T\mid F).
    \label{eq:descent-obstruction}
\end{equation}
The semantic effect descends almost surely to a predictive field if and only if this conditional variance is zero. Smooth descent additionally requires a smooth version of $\overline Z_T$ on the relevant support.
\end{theorem}

\begin{proof}
Square-integrable tangent-valued random variables form a Hilbert direct integral with inner product
\begin{equation}
    \langle U,V\rangle=\E[g_{F(H)}(U,V)].
\end{equation}
By Lemma~\ref{lem:direct-integral}, conditional expectation maps this space into the $\sigma(F)$-measurable sections and is the orthogonal projection onto them; in particular that subspace is closed. Equivalently, decompose
\begin{equation}
\begin{split}
    Z_T-Y(F)
    ={}&\bigl(Z_T-\overline Z_T(F)\bigr)\\
    &+\bigl(\overline Z_T(F)-Y(F)\bigr).
\end{split}
\end{equation}
The two terms are orthogonal after conditioning on $F$, which proves \eqref{eq:semantic-pythagoras} and \eqref{eq:descent-obstruction}.
\end{proof}

\begin{corollary}[Injective maps have no quotient obstruction]\label{cor:injective-obstruction}
If $F$ is injective on the support of $H$ and its inverse on $F(\operatorname{supp}H)$ is measurable, then
\begin{equation}
    \Var_g(Z_T\mid F)=0
    \qquad\text{almost surely}.
\end{equation}
For the affine softmax decoder, this applies whenever
\begin{equation}
    \mathcal N=\{u:Wu\in\operatorname{span}\{\bm 1\}\}=\{0\},
\end{equation}
equivalently whenever the pullback Fisher matrix \eqref{eq:softmax-metric} is positive definite.
\end{corollary}

\begin{proof}
Under the stated injectivity, $H=F^{-1}(F(H))$ almost surely, so every measurable function of $H$, including $Z_T(H)$, is measurable with respect to $\sigma(F(H))$. Its conditional variance is therefore zero. The affine-softmax equivalence follows from
\begin{equation}
    F(h)=F(h')
    \quad\Longleftrightarrow\quad
    W(h-h')\in\operatorname{span}\{\bm 1\}.
\end{equation}
When $\mathcal N=\{0\}$, Lemma~\ref{lem:affine-head-structure} also shows that
$F$ is an embedding, so its inverse on the image is continuous and hence
measurable.
\end{proof}

Corollary~\ref{cor:injective-obstruction} is a scope condition, not a defect of Theorem~\ref{thm:semantic-obstruction}. The obstruction is informative at a pre-quotient representation with nontrivial predictive fibers or at an earlier layer whose remaining network is noninjective. It is identically inert at a minimal full-rank affine head. If a smooth coarsening $S:P\to\bar P$ is declared, the theorem applies anew to $\bar F=S\circ F$, the effect $\dd\bar F\,X_T$, and the metric and tangent bundle of the coarsened image. Merely conditioning the original $P$-tangent field on $S(F)$ can mix distinct tangent spaces and is not covered by the theorem. Approximate conditioning or finite bins likewise define a different, resolution-dependent diagnostic.

\begin{corollary}[Connection-restricted decomposition]\label{cor:connection-decomposition}
Let $\mathcal C_\alpha$ be any nonempty class of square-integrable predictive
fields selected by an $\alpha$-connection, for example fields satisfying
$\nabla^{(\alpha)}Y=0$ on a specified chart. If a bounded
covariant-derivative residual is used instead, both its Sobolev domain and its
topology are part of the definition of $\mathcal C_\alpha$. Then
\begin{equation}
\begin{split}
    \inf_{Y\in\mathcal C_\alpha}\mathcal E_T(Y)
    ={}&\underbrace{\E\Var_g(Z_T\mid F)}_{\text{predictive insufficiency}}\\
    &+\underbrace{\inf_{Y\in\mathcal C_\alpha}
    \E\|\overline Z_T(F(H))-Y(F(H))\|_{g_{F(H)}}^2}_{\text{connection mismatch}}.
\end{split}
\label{eq:connection-decomposition}
\end{equation}
\end{corollary}

\begin{proof}
Equation~\eqref{eq:semantic-pythagoras} holds for every $Y\in\mathcal C_\alpha$, and its conditional-variance term is independent of $Y$. Taking the infimum over $\mathcal C_\alpha$ gives \eqref{eq:connection-decomposition}.
\end{proof}

\begin{corollary}[Cross-model semantic transfer]\label{cor:cross-model-semantic}
Let $P_A,P_B$ be predictive quotients and let $\Phi$ be a $C^1$
diffeomorphism between open neighborhoods containing the relevant supports in
$P_A$ and $P_B$. Write $Q_i=F_i(H_i)$ for the random
predictive state of model $i$.
Let $\nu_i=\mathcal L(Q_i)$ and assume
\begin{equation}
    (\Phi^{-1})_*\nu_B\ll\nu_A.
    \label{eq:cross-model-domination}
\end{equation}
Choose any $\nu_A$-almost-everywhere version of $\overline Z_A$ and define
\begin{equation}
    Y_{A\to B}(p_B)=\dd\Phi_{\Phi^{-1}(p_B)}
    \overline Z_A(\Phi^{-1}(p_B)).
\end{equation}
Assume that this transferred field is measurable and square-integrable under
$\nu_B$ in the $g_B$ norm. Then its $\nu_B$-equivalence class is independent of
the chosen $\nu_A$-version, and
\begin{equation}
\begin{split}
    \E\|Z_B-Y_{A\to B}(Q_B)\|_{g_B}^2
    ={}&\underbrace{\E\|Z_B-\overline Z_B(Q_B)\|_{g_B}^2}_{\text{model-B quotient obstruction}}\\
    &+\underbrace{\E\|\overline Z_B-\dd\Phi\,\overline Z_A\|_{g_B}^2}_{\text{cross-model field mismatch}}.
\end{split}
\label{eq:cross-model-field-decomposition}
\end{equation}
\end{corollary}

\begin{proof}
Absolute continuity in \eqref{eq:cross-model-domination} implies that two
$\nu_A$-versions of $\overline Z_A$ agree after evaluation at
$\Phi^{-1}(Q_B)$ almost surely. Thus $Y_{A\to B}$ is canonically defined as an
element of $L^2(\nu_B;TP_B,g_B)$ under the stated measurability and
integrability hypothesis.
Apply Theorem~\ref{thm:semantic-obstruction} to model B with the measurable predictive field $Y=Y_{A\to B}$. Its first term is model B's conditional Fisher variance, and its second term is
\begin{equation}
    \E\|\overline Z_B(Q_B)-Y_{A\to B}(Q_B)\|_{g_B}^2
    =\E\|\overline Z_B-\dd\Phi\,\overline Z_A\|_{g_B}^2,
\end{equation}
where the final expression evaluates $\overline Z_A$ at $\Phi^{-1}(Q_B)$. This is \eqref{eq:cross-model-field-decomposition}.
\end{proof}

Identity \eqref{eq:cross-model-field-decomposition} is measured in $g_B$ and requires only that the transferred field be well defined and measurable. If $\Phi$ is additionally a Fisher isometry, the mismatch has an intrinsic metric-preserving interpretation. Otherwise it is relative to the proposed alignment and its distortion must be reported separately.

No choice of $\alpha$ can reduce the first term of \eqref{eq:connection-decomposition}. On a genuinely noninjective predictive map, a strictly positive population conditional Fisher variance therefore proves that no predictive-state-only field, and hence no choice of $\alpha$-transport on that quotient, can represent the specified infinitesimal effect exactly. At an injective final head the term vanishes tautologically and is not an empirical diagnostic.

\section{Operational tests and training implications}

The theorems suggest a hierarchy of increasingly strong cross-model agreement tests:
\begin{equation}
\begin{array}{ccl}
\text{Level 0} &:& \text{output KL agreement},\\
\text{Level 1} &:& \text{Fisher metric agreement},\\
\text{Level 2} &:& \alpha\text{-connection and transport agreement},\\
\text{Level 3} &:& \text{semantic-field equivariance},\\
\text{Level 4} &:& \text{curvature and holonomy agreement}.
\end{array}
\end{equation}
Agreement at one level does not imply the next without the regularity and sufficiency conditions identified above.

\subsection{Transport commutation score}

For an estimated alignment $\Phi$ and a sampling law supported on paths and
nonzero tangent vectors, define
\begin{equation}
    \mathcal D_{\mathrm{PT}}^{(\alpha)}(A,B;\Phi)
    =\E_{\gamma,v}
    \frac{\|\dd\Phi_{\gamma(1)} P_\gamma^{A,\alpha}v
    -P_{\Phi\gamma}^{B,\alpha}\dd\Phi_{\gamma(0)} v\|_{g_B(\Phi(\gamma(1)))}}
    {\|\dd\Phi_{\gamma(0)}v\|_{g_B(\Phi(\gamma(0)))}}.
    \label{eq:pt-score}
\end{equation}
The numerator is measured in model $B$'s Fisher norm at the endpoint and normalized by the aligned input in model $B$'s Fisher norm at the base point. Since $\Phi$ is a local diffeomorphism, the denominator is nonzero for $v\ne0$. The integrand is invariant under separate reparameterizations of either chart. The expected score has the same invariance provided the sampling law of $(\gamma,v)$ is pushed forward with the reparameterization; a newly chosen coordinate-uniform sampling law need not be equivalent. The relative integrand does not decrease merely because every alignment derivative is multiplied by a common scalar. Because $\Phi$ also changes the compared path, alignment scale and distortion must still be audited. Define
\begin{equation}
    L(\Phi)=\sup_{x\in K}\|\dd\Phi_x\|_{g_A\to g_B},
    \qquad
    L(\Phi^{-1})=\sup_{x\in K}\|\dd\Phi_x^{-1}\|_{g_B\to g_A},
    \qquad
    \kappa(\Phi)
    =L(\Phi)L(\Phi^{-1}),
\end{equation}
and report all three quantities rather than $\kappa$ alone: uniform rescaling changes the two Lipschitz factors inversely while leaving their product unchanged. Alignment fitting must impose a preregistered scale or bi-Lipschitz range, and scores computed under different unconstrained alignments are not comparable. Theorem~\ref{thm:exact-naturality} makes this score zero under exact predictive equivalence. Corollary~\ref{cor:cross-model-commutator} bounds its aligned-norm counterpart under approximate equivalence; converting that bound to the relative Fisher score incurs the declared norm-equivalence and lower-singular-value factors. A practical study should test whether \eqref{eq:pt-score} predicts held-out steering transfer or task generalization beyond output KL and representation-only similarities such as CKA.

\subsection{Selecting the connection}

On the minimal quotient of a linear softmax head, $\alpha=1$ is ordinary
vector reuse modulo the predictive null space, $\alpha=0$ is
Fisher-metric-compatible Levi--Civita transport, and $\alpha=-1$ is mixture
transport. A predictive semantic field $Y_T$ can be evaluated through
\begin{equation}
    \mathcal R_T(\alpha)=\E\|\nabla^{(\alpha)}Y_T\|_g^2.
\end{equation}
The minimizer is interpreted relative to the tested field and chart:
\begin{itemize}[leftmargin=*]
    \item a minimum at the exponential connection means that the tested field is closest to exponential-parallel; at the affine softmax head this is constant natural-coordinate vector reuse;
    \item a minimum at the Levi-Civita connection means that the tested field is closest to the metric-compatible transport;
    \item a minimum at the mixture or an intermediate $\alpha$ identifies the corresponding affine connection as the closest member of the tested family; intermediate $\alpha$-connections need not be flat;
    \item on a verified noninjective predictive map, positive population conditional variance proves predictive insufficiency independently of the minimizing connection;
    \item on such a map, small estimated conditional variance together with large residuals for every tested $\alpha$ motivates an operation-specific connection or a nonparallel field model.
\end{itemize}

The construction of the tested field must be declared separately from the connection comparison. A finite hidden difference $h(Tc)-h(c)$ is the exponential logarithm of its endpoints at an affine head and therefore intentionally tests transfer of the conventional exponential-coordinate intervention. A connection-neutral infinitesimal test instead specifies a continuous intervention $T_t$ and uses
\begin{equation}
    X_T(h(c))=\left.\frac{\dd}{\dd t}\right|_{t=0}h(T_t c).
\end{equation}
For finite endpoints lying in a declared normal neighborhood where the relevant
affine exponential map is defined and the logarithm is single valued, each
candidate can be evaluated with its own
$\Log^{(\alpha)}$--transport--$\Exp^{(\alpha)}$ construction. Outside such a
neighborhood this notation does not define a canonical operation. Any
minimizing $\alpha\ne1$ remains a descriptive best fit within the tested
family; it does not by itself identify a mechanism implemented by optimization.

\subsection{Training regimes}

The risk-to-transport theorems indicate why derivative regularity matters. A causal training study can compare cross-entropy-only models with matched-loss models trained using
\begin{equation}
    \mathcal L_{\mathrm{total}}
    =\mathcal L_{\mathrm{CE}}
    +\lambda\widehat{\Var}_g(\dd F X_T\mid F)
    +\mu\|\nabla^{(\alpha)}\overline Z_T\|_g^2
    +\nu\mathcal R_{\mathrm{smooth}}.
    \label{eq:training-objective}
\end{equation}
The first regularizer is included only for a declared noninjective predictive map and encourages the operation to descend through that map. Under coarsening, every occurrence of $F$, $\dd F X_T$, and $g$ in this term must be replaced by $\bar F=S\circ F$, $\dd\bar F X_T$, and the coarsened-image metric. It must be omitted at an injective final head, where it is identically zero. The second selects connection-relative consistency. For Levi--Civita transport, the final term should control the square-root predictive map in the Sobolev norm used below; for general $\alpha$, it must additionally control the required score moments. Applying either theorem requires verifying the corresponding uniform bound. The empirical protocol treats behavioral improvement as the primary outcome and the geometric regularizers as diagnostic quantities.

A no-GPU pilot can use frozen small models and one- or two-dimensional soft-prompt charts. Finite differences estimate $F$, $\partial F$, and $\partial^2F$; the model's output head provides exact vocabulary probabilities. The main comparisons are output KL, metric defect, the transport score \eqref{eq:pt-score}, connection residual, and held-out intervention transfer. Conditional variance is added only at a representation with verified nontrivial predictive fibers or after a complete coarsened map, effect, and image geometry are explicitly defined.

An inferential curvature study must also separate semantic structure from Fisher-spectrum alignment. Isotropic Fisher--Haar planes answer a literal isotropic null but are not spectrum matched. A task-specific analysis should preserve the semantic plane's leverage across preregistered Fisher spectral bands, using signs for singleton bands and basis-invariant Haar rotations inside repeated or clustered bands. Fine relative-gap bands can reduce the control orbit mostly to sign flips, so the analysis must be repeated with preregistered coarser relative-log bands and conclusions that change with the band rule must be labeled control-sensitive. A plus-one control-tail rank is a calibrated randomization $p$-value only under the explicit conditional invariance null for its declared block-orthogonal group; direction and multiplicity must also be preregistered. First-order $\alpha$ fits must be extrapolated to zero edge length or replaced by integrated transport; their finite-edge parameter error is generically first order in Fisher length.

\section{Limitations}

\paragraph{Continuous charts.}
Tokenized contexts are discrete. Derivatives and connections require an explicit continuous chart, such as soft prompts, activation interventions, or a justified local manifold model.

\paragraph{Boundary and rank.}
Softmax probabilities are positive but can be extremely small. Uniform tokenwise floors are mathematically clean and practically strong, and the general-$\alpha$ probability-coordinate bound is not a useful finite-scale certificate when its constants are dominated by that floor. The square-root Levi--Civita theorem removes this explicit dependence, and affine heads admit score-moment formulas, but Fisher rank loss remains a genuine instability. In particular, the metric-compatible transport branch contains $\lambda^{-2}$ and $\lambda^{-3}$ terms after $D_{\LC}$ is expanded. The historical pilot's pointwise spectral floors worsen by more than two orders of magnitude between initialization and the final checkpoint, causing much larger inverse-power factors even before the unmeasured derivative envelopes are inserted; Appendix~\ref{app:sqrt-audit} records the values. Thus vocabulary independence is not a claim of a nonvacuous trained-model certificate. Quotienting, coarse-graining, or restriction to a preregistered transverse immersed chart is necessary for connection-level analysis near rank loss; a merely pointwise horizontal spectral subspace suffices only for pointwise metric diagnostics unless an integrable chart is also supplied. An unreported ridge changes the geometry.

\paragraph{Architectural support.}
The ambient affine-head family permits arbitrary interventions in its chosen hidden coordinates. Natural post-LayerNorm states satisfy architectural constraints and need not fill that family or remain on their support after vector addition. Claims about natural-state geometry therefore require a chart tangent to the reachable support or a pullback through the normalization map; ambient intervention geometry and support-intrinsic geometry must not be conflated.

\paragraph{Different tokenizers.}
The exact theorem assumes a common outcome space. A tokenizer map generally acts as a Markov morphism. Fisher geometry is preserved only under a statistically sufficient or reversible identification; otherwise it contracts \citep{ayJostLeSchwachhofer2015}.

\paragraph{Population versus empirical loss.}
Equation \eqref{eq:ce-decomposition} concerns population risk, realizability, and global optimization. Finite empirical loss introduces estimation and optimization errors not treated here.

\paragraph{Alignment is not identified automatically.}
The stability theorem evaluates a proposed $\Phi$. Constructing a statistically identifiable alignment from finite observations, especially across different latent dimensions, is a separate problem. One may instead work on a shared low-dimensional intervention chart.

\paragraph{Connection dependence.}
Parallel transport compares tangent vectors at different points relative to a selected connection. At the final affine softmax head, ambient addition coincides with exponential-affine transport, while Levi-Civita and mixture transport preserve different structures. The theory therefore evaluates connection choice separately across operations, layers, and tasks.

\section{Conclusion}

Idealized pretraining makes convergence of conditional output distributions natural because excess cross-entropy is expected KL divergence. It does not by itself make convergence of Fisher metrics, affine connections, or curvature natural. Geometry depends on derivatives, and a regular oscillatory counterexample separates these notions.

Under explicit smoothness and identifiability assumptions, stronger conclusions become provable. Predictive maps that converge in $C^2$ have convergent Amari $\alpha$-connections and parallel transports; $C^3$ convergence controls curvature. For Levi--Civita transport, square-root coordinates give vocabulary-independent bounds without a uniform probability floor, and integrated KL plus a uniform Hilbert-valued Sobolev bound supplies the missing derivative convergence. Nonzero-$\alpha$ connections require additional raised-third-score-moment control, for which a tokenwise probability floor is sufficient but not necessary. The semantic conditional-variance obstruction determines whether a noninjective predictive map is rich enough to support the operation; at an injective head it vanishes by construction.

If aligned predictive maps have vanishing population KL risk together with the uniform Sobolev, boundary, and rank bounds of Theorem~\ref{thm:risk-to-transport}, their $\alpha$-transport commutators vanish. Semantic-field equivariance is a separate condition: it requires small within-model quotient variance and small cross-model conditional-mean-field mismatch. These quantities localize deviations from transfer to inadequate regularity, predictive insufficiency, field mismatch, or connection mismatch.

\appendix

\section{Square-root boundary calculus and finite-scale auditing}\label{app:sqrt-audit}

\subsection{What the square-root sphere does and does not flatten}

The classical square-root map $\psi=2\sqrt p$ sends the open categorical simplex to the positive orthant of the radius-two sphere and satisfies
\begin{equation}
    \sum_a\psi_a^2=4,
    \qquad
    g=D\psi^*D\psi.
\end{equation}
The ambient inner product is Euclidean, but the induced Fisher geometry is not globally Euclidean: the full simplex has Levi--Civita sectional curvature $1/4$. Thus the sphere representation supplies a Gram factorization of the metric, not a proof of intrinsic flatness.

The metric and first-kind Levi--Civita coefficients,
\begin{equation}
    g_{ij}=\langle\partial_i\psi,\partial_j\psi\rangle,
    \qquad
    \Gamma^{\LC}_{ij,k}
    =\langle\partial_{ij}\psi,\partial_k\psi\rangle,
\end{equation}
are polynomial contractions of first and second immersion derivatives. Raised coefficients and tangent projections still contain $g^{-1}$, so the square-root representation removes explicit tokenwise denominators but does not remove Fisher rank conditioning.

For the cubic tensor, $p_a=\psi_a^2/4$ and $\partial_i p_a=(\psi_a/2)\partial_i\psi_a$ give the exact identity
\begin{equation}
    C_{ijk}
    =2\sum_a
      \frac{\partial_i\psi_a\,\partial_j\psi_a\,
            \partial_k\psi_a}{\psi_a}.
    \label{eq:sqrt-cubic-appendix}
\end{equation}
Consequently, square-root $C^2$ bounds alone do not uniformly bound $C$ near a vanishing component. Equation~\eqref{eq:sqrt-cubic-appendix} does not say that $C$ must diverge: its numerator can vanish sufficiently quickly, or structural score-moment identities can provide cancellation.

\begin{proposition}[Nonzero-$\alpha$ stability criterion]\label{prop:alpha-equivalence}
On a common chart where both pullback Fisher metrics are nondegenerate, put
$K_p=g_p^{-1}C_p$ and
$K_{\widetilde p}=g_{\widetilde p}^{-1}C_{\widetilde p}$. For every fixed
$\alpha\ne0$,
\begin{align}
    \|\Gamma_p^{(\alpha)}-\Gamma_{\widetilde p}^{(\alpha)}\|
    &\le
    \|\Gamma_p^{\LC}-\Gamma_{\widetilde p}^{\LC}\|
    +\frac{|\alpha|}{2}\|K_p-K_{\widetilde p}\|,
    \label{eq:alpha-equivalence-forward}\\
    \|K_p-K_{\widetilde p}\|
    &\le\frac{2}{|\alpha|}
    \left(
      \|\Gamma_p^{(\alpha)}-\Gamma_{\widetilde p}^{(\alpha)}\|
      +\|\Gamma_p^{\LC}-\Gamma_{\widetilde p}^{\LC}\|
    \right).
    \label{eq:alpha-equivalence-reverse}
\end{align}
Hence, once Levi--Civita stability is known, stability of one fixed nonzero-$\alpha$ connection is equivalent to stability of the raised third score moment. No tokenwise probability floor is logically necessary.
\end{proposition}

\begin{proof}
Both estimates are the triangle inequality applied in opposite directions to
\begin{equation}
    \Gamma^{(\alpha)}=\Gamma^{\LC}-\frac{\alpha}{2}K.
\end{equation}
\end{proof}

At an affine softmax head, write
\begin{equation}
    S_a(u)=(w_a-\E_pw_Y)^\top u.
\end{equation}
Then the first two derivatives of the square-root map satisfy
\begin{align}
    D\psi_a[u]&=\sqrt{p_a}\,S_a(u),\label{eq:affine-sqrt-first}\\
    D^2\psi_a[u,v]
    &=\sqrt{p_a}\left(\frac12S_a(u)S_a(v)-g(u,v)\right).
    \label{eq:affine-sqrt-second}
\end{align}
Because $\E_p[S(u)S(v)]=g(u,v)$, expansion of the squared norm gives
\begin{equation}
    \|D^2\psi[u,v]\|_{\ell_2}^2
    =\frac14\E_p[S_Y(u)^2S_Y(v)^2].
    \label{eq:b2-fourth-moment}
\end{equation}
By Cauchy--Schwarz, the associated bilinear operator norm has the exact pointwise form
\begin{equation}
    \|D^2\psi_x\|_{\op}
    =\frac12
      \left(
        \sup_{|u|\le1}\E_{p(x)}[S_Y(u)^4]
      \right)^{1/2}.
    \label{eq:b2-exact-fourth-envelope}
\end{equation}
Indeed, the mixed fourth moment is at most the geometric mean of the two pure fourth moments, while choosing $u=v$ at a maximizer attains the supremum. Thus the missing uniform $B_2$ envelope is exactly one half of the square root of the worst fourth score moment over chart-unit directions and points on the declared domain. Finite directional probes estimate this quantity but do not certify the supremum. Proposition~\ref{prop:score-moments} separately supplies the third moment needed by nonzero-$\alpha$ connections and the exact cancellation $\Gamma^{(1)}=0$ in natural affine-head coordinates.

The Bernoulli family in Proposition~\ref{prop:boundary-alpha-separation} proves that the extra condition cannot be discarded in general. More invariantly, because $g=d\theta^2$ there,
\begin{equation}
    \|\nabla^{(\alpha)}-\nabla^{\LC}\|_g
    =|\alpha\cot\theta|\longrightarrow\infty.
\end{equation}
Flat exponential or mixture coordinates do not contradict this statement: their coordinate changes become nonuniform near the boundary. The main theorems concern interior maps with no \emph{uniform} probability floor. If exact zeros or changing support are admitted, smoothness of $\psi$ must be imposed directly and the induced object is a Hellinger extension rather than an ordinary regular Fisher model on the open simplex.

\subsection{Metric-compatible transport bound}

Let $A=\nabla^{g,\LC}-\nabla^{\widetilde g,\LC}$. Up to the sign convention for $A$, the exact variation-of-connection identity is
\begin{equation}
\begin{split}
    P^{g,\LC}_{1\leftarrow0}
    -P^{\widetilde g,\LC}_{1\leftarrow0}
    =-\int_0^1
      P^{g,\LC}_{1\leftarrow t}\,
      A_{\gamma(t)}\!\left(
        \dot\gamma(t),
        P^{\widetilde g,\LC}_{t\leftarrow0}\,\cdot
      \right)\dd t.
\end{split}
\label{eq:lc-duhamel-appendix}
\end{equation}

The conditioning factors in a background Euclidean norm can be separated from the invariant transport statement. For two Riemannian metrics $g$ and $\widetilde g$, define the asymmetric mixed norm of $A=\nabla^{g,\LC}-\nabla^{\widetilde g,\LC}$ by
\begin{equation}
    \|A_x\|_{\widetilde g,\widetilde g\to g}
    =\sup_{\substack{\|u\|_{\widetilde g_x}\le1\\
                      \|v\|_{\widetilde g_x}\le1}}
       \|A_x(u,v)\|_{g_x}.
\end{equation}

\begin{proposition}[Intrinsic mixed-norm transport comparison]\label{prop:mixed-norm-transport}
Along every piecewise $C^1$ path $\gamma:[0,1]\to K$,
\begin{equation}
\begin{split}
    \|P^{g,\LC}_{1\leftarrow0}
      -P^{\widetilde g,\LC}_{1\leftarrow0}\|_{
          \widetilde g_{\gamma(0)}\to g_{\gamma(1)}}
    \le \int_0^1
       \|A_{\gamma(t)}\|_{\widetilde g,\widetilde g\to g}
       \|\dot\gamma(t)\|_{\widetilde g_{\gamma(t)}}\,\dd t.
    \label{eq:mixed-norm-transport}
\end{split}
\end{equation}
The reverse asymmetric estimate follows by interchanging $g$ and $\widetilde g$.
\end{proposition}

\begin{proof}
Apply the Duhamel identity \eqref{eq:lc-duhamel-appendix} to an input $z$ with $\|z\|_{\widetilde g_{\gamma(0)}}\le1$. The inner $\widetilde g$-transport preserves its $\widetilde g$-norm, and the outer $g$-transport preserves the $g$-norm of the output of $A$. The definition of the mixed norm bounds the remaining contraction pointwise; integration and the supremum over $z$ give \eqref{eq:mixed-norm-transport}.
\end{proof}

Equation~\eqref{eq:mixed-norm-transport} is coordinate invariant and contains no separate condition-number multiplier. It does not remove conditioning: instability of raising indices or approaching rank loss is absorbed into the mixed norm of $A$. The Euclidean certificate below is therefore useful when a shared background chart is operationally meaningful, while the mixed-norm form is the cleaner geometric comparison.

Metric compatibility and $\lambda I\preceq g\preceq B_1^2I$ imply
\begin{equation}
    \|P^{\LC}_{t\leftarrow s}\|_{\op}
    \le\frac{B_1}{\sqrt\lambda}
\end{equation}
for either geometry. Substitution into \eqref{eq:lc-duhamel-appendix} gives the polynomial branch of \eqref{eq:sqrt-transport-bound}.

For finite-scale reporting define
\begin{equation}
    \tau_\gamma=\frac{B_1B_2}{\lambda}L_\gamma,
    \qquad
    \kappa_{\mathrm{env}}=\frac{B_1^2}{\lambda}.
\end{equation}
The transport theorem reads
\begin{equation}
    \|P_\gamma^{p,\LC}-P_\gamma^{\widetilde p,\LC}\|_{\op}
    \le L_\gamma D_{\LC}
       \min\{e^{\tau_\gamma},\kappa_{\mathrm{env}}\}.
    \label{eq:finite-scale-certificate}
\end{equation}
At one point, $\|D\psi\|_{\op}^2=\lambda_{\max}(g)$, so pointwise-tight constants give $B_1^2/\lambda=\kappa(g)$. If the uniform envelopes are chosen tightly, they give
\begin{equation}
    \kappa_{\mathrm{env}}
    =\frac{\sup_K\lambda_{\max}(g)}
           {\inf_K\lambda_{\min}(g)}
    \ge\sup_K\kappa(g),
\end{equation}
and equality need not hold. For nonminimal choices of $B_1$ and $\lambda$, the
defined $\kappa_{\mathrm{env}}=B_1^2/\lambda$ is larger still. Both quantities
are relative to the declared background chart. Expanding the metric-compatible
branch makes its rank dependence explicit:
\begin{equation}
\begin{split}
    L_\gamma\frac{B_1^2}{\lambda}D_{\LC}
    =L_\gamma\left[
      \frac{B_1^2(B_1\delta_2+B_2\delta_1)}{\lambda^2}
      +\frac{2B_1^4B_2\delta_1}{\lambda^3}
    \right].
    \label{eq:expanded-polynomial-certificate}
\end{split}
\end{equation}
Consequently, removing explicit vocabulary and minimum-probability factors does not by itself make the Euclidean certificate informative near rank loss.

The historical pilot artifact \texttt{results/pythia\_14m\_pilot.json} gives the following pointwise base-state diagnostic in one fixed declared chart. The inverse powers use the smallest observed eigenvalue at each checkpoint and are not substitutes for the uniform pathwise floor required by the theorem.
\begin{center}
\small
\begin{tabular}{lrrrr}
\toprule
Checkpoint & observed $\lambda_{\min}$ range & worst $\lambda^{-1}$ & worst $\lambda^{-2}$ & worst $\lambda^{-3}$\\
\midrule
initialization & $[2.770,2.778]\times10^{-3}$ & $3.61\times10^2$ & $1.30\times10^5$ & $4.71\times10^7$\\
step143000 & $[1.199,1.391]\times10^{-5}$ & $8.34\times10^4$ & $6.96\times10^9$ & $5.80\times10^{14}$\\
\bottomrule
\end{tabular}
\end{center}
Relative to initialization, the worst displayed inverse, inverse-square, and inverse-cube factors grow by approximately $2.31\times10^2$, $5.34\times10^4$, and $1.23\times10^7$, respectively. These comparisons are chart dependent but meaningful for this fixed-chart diagnostic. They show why a trained-model certificate can be numerically vacuous even when the theorem is correct.

The pilot does not report $B_2$, $\delta_1$, $\delta_2$, or pathwise extrema, so no complete numerical value of \eqref{eq:finite-scale-certificate} is justified by the committed artifacts. The spectra motivate a conditioning audit but cannot by themselves imply a threshold such as $B_2L_\gamma\lesssim10^{-4}$ or establish that the theorem is intrinsically short-path. A numerical certificate requires measuring every term in \eqref{eq:finite-scale-certificate} along the actual comparison path and reporting whether the exponential or metric-compatible branch is active.

\section{Deferred proofs}\label{app:deferred}

This appendix proves the auxiliary identities, inequalities, and
functional-analytic facts used in the main text. Each is elementary or
classical, but several are load bearing. In particular
Lemmas~\ref{lem:hellinger-kl}, \ref{lem:hilbert-interpolation}, and
\ref{lem:hilbert-embedding} are the entire route from the training risk to
\emph{vocabulary-independent} Levi--Civita transport, so they are proved here
rather than asserted.

\subsection{The probability-coordinate connection identity}

\begin{lemma}[Lower-index $\alpha$-coefficients]\label{lem:lower-alpha}
Let $p:U\to\simplex$ be $C^2$ with $g_p$ nondegenerate. Then, in the
coordinates of \eqref{eq:pullback-metric}--\eqref{eq:pullback-cubic},
\begin{equation}
    \Gamma^{\LC}_{ij,k}
    =\sum_a\frac{\partial_{ij}p_a\,\partial_kp_a}{p_a}
     -\tfrac12 C_{ijk},
    \label{eq:lc-probability-coordinates}
\end{equation}
and consequently \eqref{eq:lower-alpha} holds for every $\alpha\in\R$.
\end{lemma}

\begin{proof}
Differentiating \eqref{eq:pullback-metric},
\begin{equation}
    \partial_kg_{ij}
    =\sum_a\left[
      \frac{\partial_{ik}p_a\,\partial_jp_a}{p_a}
      +\frac{\partial_ip_a\,\partial_{jk}p_a}{p_a}
      -\frac{\partial_ip_a\,\partial_jp_a\,\partial_kp_a}{p_a^2}
    \right],
    \label{eq:metric-derivative-expansion}
\end{equation}
and cyclically for $\partial_ig_{jk}$ and $\partial_jg_{ik}$. Collect the
Koszul combination $\partial_ig_{jk}+\partial_jg_{ik}-\partial_kg_{ij}$ term by
term. The term $\partial_{ij}p_a\,\partial_kp_a/p_a$ occurs once in each of the
first two expansions and not in the third, contributing a factor two. The term
$\partial_{ik}p_a\,\partial_jp_a/p_a$ occurs in the first with sign $+$ and in
the third with sign $-$, and cancels; likewise
$\partial_ip_a\,\partial_{jk}p_a/p_a$ cancels between the second and the third.
The cubic terms contribute $(-1)+(-1)-(-1)=-1$ times $C_{ijk}$. Hence
\begin{equation}
    \partial_ig_{jk}+\partial_jg_{ik}-\partial_kg_{ij}
    =2\sum_a\frac{\partial_{ij}p_a\,\partial_kp_a}{p_a}
     -C_{ijk},
\end{equation}
and halving gives \eqref{eq:lc-probability-coordinates}. Lowering the final
index of \eqref{eq:alpha-connection} with $g_p$ gives
$\Gamma^{(\alpha)}_{ij,k}=\Gamma^{\LC}_{ij,k}-\tfrac{\alpha}{2}C_{ijk}$, which
together with \eqref{eq:lc-probability-coordinates} is \eqref{eq:lower-alpha}.
\end{proof}

\begin{lemma}[Duality of $\pm\alpha$]\label{lem:alpha-duality}
For every $\alpha\in\R$,
$\partial_kg_{ij}=\Gamma^{(\alpha)}_{ki,j}+\Gamma^{(-\alpha)}_{kj,i}$.
\end{lemma}

\begin{proof}
By \eqref{eq:lower-alpha} and total symmetry of $C$,
\begin{equation}
\begin{split}
    \Gamma^{(\alpha)}_{ki,j}+\Gamma^{(-\alpha)}_{kj,i}
    ={}&\sum_a\frac{\partial_{ki}p_a\,\partial_jp_a
        +\partial_{kj}p_a\,\partial_ip_a}{p_a}\\
      &-\left(\tfrac{1+\alpha}{2}+\tfrac{1-\alpha}{2}\right)C_{ijk},
\end{split}
\end{equation}
and the right-hand side is exactly
\eqref{eq:metric-derivative-expansion}.
\end{proof}

\begin{lemma}[Identification of $\alpha=\pm1$]\label{lem:alpha-identification}
Let $p_\theta(a)=\exp(\theta\cdot T(a)-\psi(\theta))$ be a regular exponential
family in minimal natural coordinates $\theta$, and put
$\eta_i=\partial_i\psi$. Then
\begin{equation}
    \Gamma^{(\alpha)}_{ij,k}=\tfrac{1-\alpha}{2}C_{ijk},
    \qquad
    g_{ij}=\partial_{ij}\psi,
    \qquad
    C_{ijk}=\partial_{ijk}\psi .
    \label{eq:exponential-family-alpha}
\end{equation}
In particular $\nabla^{(1)}$ has vanishing coefficients in $\theta$, so
$\theta$ is $\nabla^{(1)}$-affine and constant coordinate vectors are
$\nabla^{(1)}$-parallel; and $\eta$ is $\nabla^{(-1)}$-affine, so
$\nabla^{(-1)}$ is the mixture connection.
\end{lemma}

\begin{proof}
Here $\log p_\theta(a)=\theta\cdot T(a)-\psi(\theta)$, so
$S_i(a)=T_i(a)-\partial_i\psi$ and $H_{ij}(a)=-\partial_{ij}\psi$ is
independent of $a$. Because $\E_p[S_k]=0$, the first term of
\eqref{eq:score-alpha} vanishes, which gives the first identity of
\eqref{eq:exponential-family-alpha}. The remaining two are the standard
cumulant identities $\E_p[S_iS_j]=\partial_{ij}\psi$ and
$\E_p[S_iS_jS_k]=\partial_{ijk}\psi$ for the cumulant generating function
$\psi$. Setting $\alpha=1$ gives $\Gamma^{(1)}\equiv0$, proving the first
claim.

For $\alpha=-1$, Lemma~\ref{lem:alpha-duality} with $\Gamma^{(1)}\equiv0$ gives
$\Gamma^{(-1)}_{kj,i}=\partial_kg_{ij}=\partial_{ijk}\psi$, hence
$\Gamma^{(-1)}{}^l{}_{ij}=g^{lk}\partial_{ijk}\psi$. The coordinate change
$\theta\mapsto\eta=\nabla\psi(\theta)$ has Jacobian
$\partial\eta_i/\partial\theta^j=g_{ij}$, so
$\partial/\partial\eta^a=g^{ai}\partial/\partial\theta^i$ and
\begin{equation}
    \nabla^{(-1)}_{\partial^\eta_a}\partial^\eta_b
    =g^{ai}\Bigl[
        (\partial_ig^{bj})\,\partial_j
        +g^{bj}g^{lk}\partial_{ijk}\psi\,\partial_l
      \Bigr].
\end{equation}
Differentiating $g^{bj}g_{jq}=\delta^b_q$ gives
$\partial_ig^{bj}=-g^{bp}g^{qj}\partial_{ipq}\psi$, so the first summand is
$-g^{ai}g^{bp}g^{qj}\partial_{ipq}\psi\,\partial_j$. Renaming $j\to p$,
$l\to j$, $k\to q$ in the second summand turns it into
$+g^{ai}g^{bp}g^{jq}\partial_{ipq}\psi\,\partial_j$. The two cancel, so
$\nabla^{(-1)}_{\partial^\eta_a}\partial^\eta_b=0$ and $\eta$ is
$\nabla^{(-1)}$-affine.
\end{proof}

\subsection{Elementary inequalities}

\begin{lemma}[Hellinger--Kullback]\label{lem:hellinger-kl}
For $p,q\in\simplex$,
$\sum_a(\sqrt{p_a}-\sqrt{q_a})^2\le\KL(p\Vert q)$.
\end{lemma}

\begin{proof}
Let $\rho=\sum_a\sqrt{p_aq_a}$ be the Bhattacharyya coefficient, so that
$\sum_a(\sqrt{p_a}-\sqrt{q_a})^2=2(1-\rho)$. By concavity of the logarithm and
Jensen's inequality under $p$,
\begin{equation}
    -\tfrac12\KL(p\Vert q)
    =\sum_ap_a\log\sqrt{\frac{q_a}{p_a}}
    \le\log\sum_ap_a\sqrt{\frac{q_a}{p_a}}
    =\log\rho ,
\end{equation}
that is, $\KL(p\Vert q)\ge-2\log\rho$. Since $\log t\le t-1$ for $t>0$ we have
$-2\log\rho\ge2(1-\rho)$, and combining the two displays proves the claim.
\end{proof}

\begin{lemma}[Two-model $L^1$ bound]\label{lem:two-model-l1}
Inequality \eqref{eq:two-model-pinsker} holds.
\end{lemma}

\begin{proof}
Pointwise in $c$, the triangle inequality in $L^1$ and Pinsker's inequality
give $\|q_A-q_B\|_1\le\sqrt{2\KL(P\Vert q_A)}+\sqrt{2\KL(P\Vert q_B)}$. Take
expectations over $C$ and apply Jensen's inequality to the concave function
$t\mapsto\sqrt t$, which replaces $\E_C\sqrt{2\KL(P\Vert q_i)}$ by at most
$\sqrt{2\,\E_C\KL(P\Vert q_i)}=\sqrt{2\epsilon_i}$. The Jensen step is the
explicit argument that converts the intermediate quantity
$\E_C\sqrt{\KL}$ into a function of the excess risk $\E_C\KL$.
\end{proof}

\subsection{Structure of the affine softmax head}

\begin{lemma}[Kernel and embedding]\label{lem:affine-head-structure}
Let $F(h)=\operatorname{softmax}(Wh+b)$ with $W\in\R^{N\times d}$, and let
$\mathcal N=\{u:Wu\in\operatorname{span}\{\bm 1\}\}$. Then
\begin{enumerate}[leftmargin=*,itemsep=0.2em]
    \item $u^\top G(h)u=\Var_{p_h}(w_Y^\top u)$ for every $u$, and hence
      $\ker G(h)=\mathcal N$ for every $h$;
    \item $F(h)=F(h')$ if and only if $h-h'\in\mathcal N$;
    \item the induced map $\bar F:\R^d/\mathcal N\to\simplex$ is a smooth
      embedding.
\end{enumerate}
\end{lemma}

\begin{proof}
(i) Expanding \eqref{eq:softmax-metric},
\begin{equation}
    u^\top G(h)u
    =\sum_ap_a(w_a^\top u)^2-\Bigl(\sum_ap_aw_a^\top u\Bigr)^2
    =\Var_{p_h}(w_Y^\top u),
\end{equation}
which vanishes if and only if $a\mapsto w_a^\top u$ is constant on
$\{a:p_a>0\}$. Since $p_h$ is strictly positive this is exactly
$Wu\in\operatorname{span}\{\bm 1\}$.

(ii) The softmax map is injective modulo additive constants on logits, so
$F(h)=F(h')$ if and only if $W(h-h')\in\operatorname{span}\{\bm 1\}$.

(iii) Fix a reference outcome $a_0$ and let $L:\simplex\to\R^{N-1}$ have
components $L(p)_a=\log(p_a/p_{a_0})$ for $a\ne a_0$. Then $L$ is a
diffeomorphism, its inverse being the softmax that pins coordinate $a_0$ to
zero. Moreover $L\circ F$ is the affine map
\begin{equation}
    h\longmapsto
    \bigl((w_a-w_{a_0})^\top h+(b_a-b_{a_0})\bigr)_{a\ne a_0},
\end{equation}
whose linear part has kernel exactly $\mathcal N$ by (ii). Hence
$L\circ\bar F$ is an injective affine map of the finite-dimensional vector
space $\R^d/\mathcal N$ onto a closed affine subspace of $\R^{N-1}$; such a map
is a homeomorphism onto its image and an immersion. Composing with the
diffeomorphism $L^{-1}$ shows that $\bar F$ is an injective immersion that is a
homeomorphism onto its image, that is, a smooth embedding. (Injective
immersions need not be embeddings in general, which is why the intermediate
log-odds chart is used.)
\end{proof}

\begin{remark}[Quotients in the general case]\label{rem:quotient}
Lemma~\ref{lem:affine-head-structure} is elementary because $\mathcal N$ is a
fixed linear subspace. For a general smooth $F:M\to\simplex$ we do not infer a
manifold structure on $M/\!\sim_F$ from regular-looking fibers alone. Whenever
a global predictive quotient is used, the factorization
$F=\bar F\circ\pi$ declared in Section~3 is an assumption: $\pi:M\to M_F$ is a
surjective submersion with exactly the predictive fibers and
$\bar F:M_F\to\simplex$ is an embedding. A standard quotient-manifold theorem
would alternatively require the equivalence-relation graph to be a closed
embedded submanifold of $M\times M$ and either coordinate projection of that
graph to be a submersion; we do not invoke that theorem below. Without a global
quotient, every estimate is applied on an explicitly declared transverse
immersed chart. This avoids assigning an ordinary Levi--Civita connection to a
possibly nonintegrable horizontal distribution.
\end{remark}

\begin{lemma}[Pullback nondegeneracy]\label{lem:pullback-nondegenerate}
If $F:M\to\simplex$ is a smooth immersion then $g_F$ is a Riemannian metric.
\end{lemma}

\begin{proof}
Since $\sum_ap_a\equiv1$ we have $\sum_a\partial_ip_a=0$, so $\dd F$ takes
values in $T_0=\{u\in\R^N:\sum_au_a=0\}$. On $T_0$ the form
$g^\Delta(u,u)=\sum_au_a^2/p_a$ is positive definite. Hence
$g_F(u,u)=g^\Delta(\dd Fu,\dd Fu)=0$ forces $\dd Fu=0$, and $u=0$ because $F$
is an immersion.
\end{proof}

\subsection{The two quantitative steps in Theorem~\ref{thm:c2-stability}}

\begin{lemma}[Lipschitz dependence on the regularity class]\label{lem:lipschitz-class}
Fix $\eta,B_2>0$ and set
\begin{equation}
    \mathcal J
    =\Bigl\{(P,P',P''):
      P_a\ge\eta,\ \textstyle\sum_aP_a=1,\
      |P'|,|P''|\le B_2\Bigr\}.
\end{equation}
The map sending $(P,P',P'')$ to the tuple $(g,Dg,C)$ assembled from
\eqref{eq:pullback-metric}, \eqref{eq:pullback-cubic}, and
\eqref{eq:metric-derivative-expansion} is Lipschitz on $\mathcal J$ with a
constant depending only on $m,N,\eta^{-1},B_2$.
\end{lemma}

\begin{proof}
$\mathcal J$ is an intersection of half-spaces, one hyperplane, and closed norm
balls, hence convex and compact. Every component of $(g,Dg,C)$ is a rational
function of $(P,P',P'')$ whose only denominators are powers of $P_a$, so the
map is $C^\infty$ on the open set $\{P_a>\eta/2\}\supset\mathcal J$ and its
differential is bounded on the compact set $\mathcal J$ by some
$L(m,N,\eta^{-1},B_2)$. Convexity guarantees that the segment joining any two
points of $\mathcal J$ stays inside $\mathcal J$, so the mean-value inequality
may be applied along that segment and yields the Lipschitz bound. Convexity is
not decorative here: the mean-value inequality fails on a general domain.
\end{proof}

\begin{lemma}[Duhamel with a cancelling growth factor]\label{lem:duhamel-cancellation}
Let $\gamma:[0,1]\to K$ be piecewise $C^1$ of Euclidean length $L_\gamma$, let
$\Gamma,\widetilde\Gamma$ be continuous connection coefficients with
$\|\Gamma\|_{C^0},\|\widetilde\Gamma\|_{C^0}\le M$, and let
$P_\gamma,\widetilde P_\gamma$ denote the induced transports. Then
\begin{equation}
    \|P_\gamma-\widetilde P_\gamma\|_{\op}
    \le L_\gamma\,e^{ML_\gamma}\,
      \|\Gamma-\widetilde\Gamma\|_{C^0}.
\end{equation}
\end{lemma}

\begin{proof}
Write $s(t)=\int_0^t|\dot\gamma|$, so $s(0)=0$ and $s(1)=L_\gamma$. Gronwall's
inequality applied to $\dot V=-\Gamma[\dot\gamma]V$ gives
$\|P_{1\leftarrow t}\|_{\op}\le\exp\bigl(M(L_\gamma-s(t))\bigr)$ and
$\|\widetilde P_{t\leftarrow0}\|_{\op}\le\exp\bigl(Ms(t)\bigr)$. Duhamel's
identity writes $P_\gamma-\widetilde P_\gamma$ as
\begin{equation}
    -\int_0^1P_{1\leftarrow t}\,
      (\Gamma-\widetilde\Gamma)[\dot\gamma(t)]\,
      \widetilde P_{t\leftarrow0}\,\dd t,
\end{equation}
whence
\begin{equation}
    \|P_\gamma-\widetilde P_\gamma\|_{\op}
    \le\int_0^1
      e^{M(L_\gamma-s(t))}e^{Ms(t)}
      \|\Gamma-\widetilde\Gamma\|_{C^0}\,|\dot\gamma(t)|\,\dd t .
\end{equation}
The two exponentials multiply to $e^{ML_\gamma}$, which does not depend on $t$
and therefore leaves the integral; the remaining $\int_0^1|\dot\gamma|$ equals
$L_\gamma$. Bounding each propagator separately by $e^{ML_\gamma}$ would give
the weaker factor $e^{2ML_\gamma}$, so this cancellation is what produces the
constant stated in \eqref{eq:pt-stability}. If the two coefficient fields carry
different bounds $M$ and $\widetilde M$, the same computation gives
$e^{\max\{M,\widetilde M\}L_\gamma}$.
\end{proof}

\subsection{Hilbert-valued Sobolev estimates with dimension-free constants}

Theorem~\ref{thm:sqrt-risk-to-transport} claims constants independent of the
vocabulary size $N$. Since $\psi_p$ takes values in $\ell_2^N$, we record
Sobolev interpolation and embedding in Hilbert-space-valued form and track
their constants. For finite $N$ the same dimension-free estimate can also be
obtained componentwise by squaring the scalar estimates and summing; a
$\sqrt N$ loss arises only from the unnecessarily crude step of bounding every
component by the full vector-valued norm before summing. Throughout,
$\mathbb H$ is a separable Hilbert space,
integrals of $\mathbb H$-valued functions are Bochner integrals, and for $s\ge0$
\begin{equation}
    \|f\|_{H^s(\R^m;\mathbb H)}^2
    =\int_{\R^m}(1+|\xi|^2)^s\|\hat f(\xi)\|_{\mathbb H}^2\,\dd\xi .
\end{equation}

\begin{lemma}[Vector-valued Plancherel]\label{lem:vector-plancherel}
The Fourier transform extends to a unitary operator of $L^2(\R^m;\mathbb H)$, with no
dependence on $\dim\mathbb H$.
\end{lemma}

\begin{proof}
For a real Hilbert space, take its complexification; the canonical embedding is
isometric, so it suffices to prove the complex case.
Fix an orthonormal basis $(e_n)$ of $\mathbb H$ and write $f=\sum_nf_ne_n$ with
$f_n=\langle f,e_n\rangle_{\mathbb H}\in L^2(\R^m)$. The Fourier transform acts
coordinatewise, $\langle\hat f,e_n\rangle=\widehat{f_n}$, and the scalar
Plancherel theorem gives $\|f_n\|_{L^2}=\|\widehat{f_n}\|_{L^2}$. By monotone
convergence,
\[
\|f\|_{L^2(\mathbb H)}^2=\sum_n\|f_n\|_{L^2}^2
=\sum_n\|\widehat{f_n}\|_{L^2}^2=\|\hat f\|_{L^2(\mathbb H)}^2.
\]
\end{proof}

\begin{lemma}[Interpolation with constant one]\label{lem:hilbert-interpolation}
For $0\le r\le s$ and $f\in H^s(\R^m;\mathbb H)$,
\begin{equation}
    \|f\|_{H^r(\R^m;\mathbb H)}
    \le\|f\|_{L^2(\R^m;\mathbb H)}^{1-r/s}
       \|f\|_{H^s(\R^m;\mathbb H)}^{r/s}.
\end{equation}
\end{lemma}

\begin{proof}
The endpoint cases $r=0$ and $r=s$ are identities, so assume $0<r<s$.
Apply H\"older's inequality with conjugate exponents $s/r$ and $s/(s-r)$ to the
factorization
\begin{equation}
    (1+|\xi|^2)^r\|\hat f(\xi)\|_{\mathbb H}^2
    =\Bigl[(1+|\xi|^2)^s\|\hat f(\xi)\|_{\mathbb H}^2\Bigr]^{r/s}
     \Bigl[\|\hat f(\xi)\|_{\mathbb H}^2\Bigr]^{(s-r)/s},
\end{equation}
which gives
$\|f\|_{H^r}^2\le\|f\|_{H^s}^{2r/s}\|f\|_{L^2}^{2(s-r)/s}$; take square roots.
The only property of $\mathbb H$ used is that $\xi\mapsto\|\hat f(\xi)\|_{\mathbb H}$ is a
nonnegative scalar function, so the constant is one for every $\mathbb H$.
\end{proof}

\begin{lemma}[Embedding with a dimension-free constant]\label{lem:hilbert-embedding}
Let $k\in\mathbb N_0$ and $r>k+m/2$. Put
\begin{equation}
    \kappa=(2\pi)^{-m/2}\max_{|\beta|\le k}
      \left(\int_{\R^m}
        \frac{|\xi|^{2|\beta|}}{(1+|\xi|^2)^{r}}\,\dd\xi
      \right)^{1/2}<\infty .
    \label{eq:embedding-constant}
\end{equation}
Then for every separable Hilbert space $\mathbb H$ and every $f\in H^r(\R^m;\mathbb H)$, the
derivatives $D^\beta f$ with $|\beta|\le k$ have continuous bounded
representatives and
\begin{equation}
    \max_{|\beta|\le k}\sup_{x\in\R^m}\|D^\beta f(x)\|_{\mathbb H}
    \le\kappa\,\|f\|_{H^r(\R^m;\mathbb H)} .
\end{equation}
\end{lemma}

\begin{proof}
Let first $f$ be an $\mathbb H$-valued Schwartz function. Fourier inversion gives the
Bochner integral
$D^\beta f(x)=(2\pi)^{-m/2}\int(i\xi)^\beta\hat f(\xi)e^{ix\cdot\xi}\,\dd\xi$.
The triangle inequality for Bochner integrals together with
$|\xi^\beta|\le|\xi|^{|\beta|}$ gives
\begin{equation}
    \|D^\beta f(x)\|_{\mathbb H}
    \le(2\pi)^{-m/2}\!\!\int_{\R^m}\!
      \frac{|\xi|^{|\beta|}}{(1+|\xi|^2)^{r/2}}\cdot
      (1+|\xi|^2)^{r/2}\|\hat f(\xi)\|_{\mathbb H}\,\dd\xi .
\end{equation}
Both factors under the integral are nonnegative \emph{scalar} functions of
$\xi$, so the Cauchy--Schwarz inequality in $L^2(\dd\xi)$ applies and yields
$\|D^\beta f(x)\|_{\mathbb H}\le\kappa_\beta\|f\|_{H^r(\R^m;\mathbb H)}$, where $\kappa_\beta$ is
the corresponding entry of \eqref{eq:embedding-constant}. In polar coordinates
the defining integral converges exactly when $2|\beta|-2r+m<0$, that is when
$r>|\beta|+m/2$, which holds for every $|\beta|\le k$ by hypothesis. The bound
is uniform in $x$ and the integrand is dominated independently of $x$, so
dominated convergence makes each $D^\beta f$ continuous. Since $\mathbb H$-valued
Schwartz functions are dense in $H^r(\R^m;\mathbb H)$ and the estimate is continuous in
the $H^r$ norm, it extends to all of $H^r(\R^m;\mathbb H)$, the limit supplying the
continuous representative.

No step refers to $\dim\mathbb H$. The Hilbert structure enters only through the
Bochner triangle inequality and Lemma~\ref{lem:vector-plancherel}, while
$\kappa$ is an integral of a scalar function over $\R^m$. For finite $N$, the
same conclusion follows from the scalar estimate by
\[
  \sum_{a=1}^N|D^\beta f_a(x)|^2
  \le \kappa_\beta^2\sum_{a=1}^N\|f_a\|_{H^r}^2;
\]
this is dimension free for exactly the same Hilbertian reason.
\end{proof}

\begin{lemma}[Transfer to domains and compact manifolds]\label{lem:hilbert-domain}
Let $\Omega\subset\R^m$ satisfy the common extension-operator hypothesis
\eqref{eq:total-extension-definition} through order $s$. Then
Lemmas~\ref{lem:hilbert-interpolation} and \ref{lem:hilbert-embedding} hold on
$\Omega$ with constants $C(\Omega,m,k,r,s)$ independent of $\mathbb H$.
They also hold on a compact $m$-dimensional Riemannian manifold for the spectral
Sobolev scale declared in Section~\ref{sec:risk-to-transport}, with constants
depending on the background manifold but not on $\mathbb H$.
\end{lemma}

\begin{proof}
On the domain, let $E=E_\Omega$ be the single scalar extension operator from
\eqref{eq:total-extension-definition}. Since all spaces here are Hilbert
spaces, the $\mathbb H$-valued extension is the Hilbert tensor product
$E\otimes I_{\mathbb H}$, and
\begin{equation}
  \|E\otimes I_{\mathbb H}\|_{H^t(\Omega;\mathbb H)\to
  H^t(\R^m;\mathbb H)}=\|E\|_{H^t(\Omega)\to H^t(\R^m)}=:C_t.
\end{equation}
Indeed, the equality is immediate first on finite orthogonal sums and then by
density; the reverse inequality follows by tensoring a scalar near-maximizer
with a unit vector in $\mathbb H$. For interpolation,
\begin{equation}
\begin{split}
    \|f\|_{H^r(\Omega;\mathbb H)}
    &\le\|Ef\|_{H^r(\R^m;\mathbb H)}
     \le\|Ef\|_{L^2}^{1-r/s}\|Ef\|_{H^s}^{r/s}\\
    &\le C_0^{1-r/s}C_s^{r/s}
      \|f\|_{L^2(\Omega;\mathbb H)}^{1-r/s}
      \|f\|_{H^s(\Omega;\mathbb H)}^{r/s},
\end{split}
\end{equation}
and for the embedding
$\|f\|_{C^k(\overline\Omega;\mathbb H)}\le\|Ef\|_{C^k(\R^m;\mathbb H)}
\le\kappa C_r\|f\|_{H^r(\Omega;\mathbb H)}$.

On a compact manifold let $L=I+\Delta_h$. By the spectral theorem, for each
$f$ there is a finite positive scalar measure $\mu_f$ on $[1,\infty)$ such that
\begin{equation}
  \|f\|_{H^t(\Omega;\mathbb H)}^2
  =\int_{[1,\infty)}\lambda^t\,\dd\mu_f(\lambda).
\end{equation}
H\"older's inequality applied to this scalar integral proves interpolation with
constant one, exactly as in Lemma~\ref{lem:hilbert-interpolation}.

For embedding, the standard equivalence between the spectral Sobolev norm and
the norm obtained from a finite atlas and subordinate partition of unity
\citep{hebey1996} reduces the scalar Sobolev theorem to
Lemma~\ref{lem:hilbert-embedding} on finitely many Euclidean chart domains.
Consequently, for every $0\le j\le k$ and $x\in\Omega$,
the scalar covariant-derivative evaluation map
\begin{equation}
  T_{x,j}:H^r(\Omega)\longrightarrow (T_x^*\Omega)^{\otimes j},
  \qquad u\longmapsto(\nabla^0)^ju(x),
\end{equation}
has operator norm at most one uniform constant $C_{r,k,\Omega}$. The
$\mathbb H$-valued evaluation is $T_{x,j}\otimes I_{\mathbb H}$ and therefore
has exactly the same operator norm. Taking the maximum over $j$ and supremum
over $x$ gives the desired dimension-free $C^k$ estimate. Smooth
$\mathbb H$-valued functions are dense in the spectral Sobolev space, and the
uniform estimate makes their derivatives converge uniformly, supplying the
continuous representatives.
\end{proof}

\begin{remark}[Where vocabulary-independence comes from]\label{rem:where-N-free}
In Theorem~\ref{thm:sqrt-risk-to-transport} the target is
$\mathbb H=\ell_2^N$.
Lemma~\ref{lem:hellinger-kl} converts the chart-integrated risk into an
$L^2(\Omega;\ell_2^N)$ bound with the absolute constant $4$;
Lemmas~\ref{lem:hilbert-interpolation}--\ref{lem:hilbert-domain} convert that
into a $C^2(\Omega;\ell_2^N)$ bound with constants depending only on
$\Omega,m,r,s$; and Theorem~\ref{thm:sqrt-lc-stability}, or its covariant form
Lemma~\ref{lem:manifold-sqrt-stability}, consumes only
$\delta_1,\delta_2,B_1,B_2,\lambda$. No step introduces a factor of $N$ or of
$\min_ap_a$. What survives is the dependence on the uniform envelope $B$ and
the spectral floor $\lambda$, whose deterioration with model scale is the
subject of Appendix~\ref{app:sqrt-audit}.
\end{remark}

\subsection{Background-covariant manifold estimates}

Let $(\Omega,h)$ be a compact Riemannian manifold and write $\nabla^0$ for the
Levi--Civita connection of the fixed background metric $h$. Covariant $C^k$
norms below use $\nabla^0$ and $h$. This formulation avoids comparing
non-tensorial Christoffel symbols across chart changes.

\begin{lemma}[Covariant probability-coordinate stability]\label{lem:chart-transfer}
Let $p,\widetilde p:\Omega\to\simplex$ be $C^2$ and suppose, for constants
$\eta,B_2,\lambda>0$,
\begin{equation}
  p_a,\widetilde p_a\ge\eta,\qquad
  \|p\|_{C^2_h}+\|\widetilde p\|_{C^2_h}\le B_2,\qquad
  g_p,g_{\widetilde p}\succeq\lambda h.
\end{equation}
For $|\alpha|\le A$, put
$A_p^{(\alpha)}=\nabla_p^{(\alpha)}-\nabla^0$ and define
$A_{\widetilde p}^{(\alpha)}$ similarly. Then
\begin{equation}
\begin{split}
  &\|g_p-g_{\widetilde p}\|_{C^1_h}
   +\|C_p-C_{\widetilde p}\|_{C^0_h}
   +\|A_p^{(\alpha)}-A_{\widetilde p}^{(\alpha)}\|_{C^0_h}\\
  &\hspace{7em}\le
  C_h\|p-\widetilde p\|_{C^2_h},
  \label{eq:covariant-probability-stability}
\end{split}
\end{equation}
where $C_h$ depends only on
$m,N,A,\eta^{-1},\lambda^{-1},B_2$ and the fixed background norm conventions.
The same data give a finite bound
\begin{equation}
  M_h=\max\{\|A_p^{(\alpha)}\|_{C^0_h},
             \|A_{\widetilde p}^{(\alpha)}\|_{C^0_h}\}.
\end{equation}
For every piecewise $C^1$ path of background length $L_\gamma$,
\begin{equation}
  \|P_\gamma^{p,\alpha}-P_\gamma^{\widetilde p,\alpha}\|_{h\to h}
  \le L_\gamma e^{M_hL_\gamma}C_h
       \|p-\widetilde p\|_{C^2_h}.
  \label{eq:covariant-probability-transport}
\end{equation}
If in addition $p,\widetilde p$ are $C^3$ with
$\|p\|_{C^3_h}+\|\widetilde p\|_{C^3_h}\le B_3$, then
\begin{equation}
  \|R_p^{(\alpha)}-R_{\widetilde p}^{(\alpha)}\|_{C^0_h}
  \le C_{R,h}\|p-\widetilde p\|_{C^3_h},
  \label{eq:covariant-curvature-stability}
\end{equation}
with $C_{R,h}$ depending additionally on $B_3$. For $\alpha=0$, sectional
curvatures converge on corresponding planes whose Fisher Gram determinants
have a common positive lower bound.
\end{lemma}

\begin{proof}
The difference of two connections is a tensor. The Koszul difference formula,
using that $\nabla^0$ is torsion free, gives the following global version of
\eqref{eq:lower-alpha}: for all tangent vectors $X,Y,Z$,
\begin{equation}
\begin{split}
  g_p\bigl(A_p^{(\alpha)}(X,Y),Z\bigr)
  =\sum_a\Bigg[
     &\frac{(\nabla^0\dd p_a)(X,Y)\,\dd p_a(Z)}{p_a}\\
     &-\frac{1+\alpha}{2}
      \frac{\dd p_a(X)\dd p_a(Y)\dd p_a(Z)}{p_a^2}
  \Bigg].
  \label{eq:covariant-lower-alpha}
\end{split}
\end{equation}
Indeed,
$2g_p(A_p^{\LC}(X,Y),Z)$ equals
$(\nabla^0_Xg_p)(Y,Z)+(\nabla^0_Yg_p)(X,Z)
-(\nabla^0_Zg_p)(X,Y)$; expanding
$g_p=\sum_a p_a^{-1}\dd p_a\otimes\dd p_a$ cancels the two mixed Hessian
terms and yields the case $\alpha=0$, after which
$A_p^{(\alpha)}=A_p^{\LC}-(\alpha/2)g_p^{-1}C_p$ gives
\eqref{eq:covariant-lower-alpha}.

The tensors $g_p$, $\nabla^0g_p$, $C_p$, and the lower-index tensor in
\eqref{eq:covariant-lower-alpha} are finite sums of smooth rational functions
of $(p,\dd p,\nabla^0\dd p)$ whose only denominators are powers of $p_a$.
The convex jet argument of Lemma~\ref{lem:lipschitz-class}, now in
background-orthonormal frames, controls their differences by
$C\|p-\widetilde p\|_{C^2_h}$. Raising the final index uses only the endpoint
identity
\[
  g_p^{-1}-g_{\widetilde p}^{-1}
  =g_p^{-1}(g_{\widetilde p}-g_p)g_{\widetilde p}^{-1}
\]
and the two bounds by $\lambda^{-1}$. This proves
\eqref{eq:covariant-probability-stability}; the same formulas without
subtraction give $M_h<\infty$ with the stated dependence.

Along $\gamma$, choose an $h$-orthonormal $\nabla^0$-parallel frame. In this
frame $\nabla_p^{(\alpha)}$-transport solves
$\dot v=-a_p(t)v$, where
$\|a_p(t)\|\le M_h\|\dot\gamma(t)\|_h$; the coefficient defect is at most
$C_h\|p-\widetilde p\|_{C^2_h}\|\dot\gamma(t)\|_h$. The frame identifies the
endpoint background norms isometrically, so Lemma~\ref{lem:duhamel-cancellation}
applies verbatim with background arclength and proves
\eqref{eq:covariant-probability-transport}. No chart subdivision occurs.

Under the $C^3$ envelope, covariantly differentiating
\eqref{eq:covariant-lower-alpha} and using
\[
  \nabla^0(g_p^{-1})
  =-g_p^{-1}(\nabla^0g_p)g_p^{-1}
\]
gives
$\|\nabla^0A_p^{(\alpha)}-\nabla^0A_{\widetilde p}^{(\alpha)}\|_{C^0_h}
\le C\|p-\widetilde p\|_{C^3_h}$; expanding the difference of the displayed
inverse identity uses only endpoint inverse bounds, exactly as in the proof of
Theorem~\ref{thm:c3-curvature}. The curvature difference follows from
\begin{equation}
\begin{split}
 R_p^{(\alpha)}(X,Y)Z
 ={}&R^0(X,Y)Z
   +(\nabla^0_XA_p^{(\alpha)})(Y,Z)
   -(\nabla^0_YA_p^{(\alpha)})(X,Z)\\
  &+A_p^{(\alpha)}(X,A_p^{(\alpha)}(Y,Z))
   -A_p^{(\alpha)}(Y,A_p^{(\alpha)}(X,Z)).
\end{split}
\end{equation}
Subtracting the two formulas cancels $R^0$ and proves
\eqref{eq:covariant-curvature-stability}. The sectional-curvature conclusion
uses background-orthonormal plane representatives and the Gram floor, as in
Theorem~\ref{thm:c3-curvature}.
\end{proof}

\begin{lemma}[Covariant square-root stability]\label{lem:manifold-sqrt-stability}
Let $\psi,\widetilde\psi:\Omega\to\mathbb H$ be $C^2$ immersions into a
Hilbert space and set $g=D\psi^*D\psi$ and
$\widetilde g=D\widetilde\psi^*D\widetilde\psi$. Suppose
\begin{equation}
  \|D\psi\|,\|D\widetilde\psi\|\le B_1,\qquad
  \|\nabla^0D\psi\|,\|\nabla^0D\widetilde\psi\|\le B_2,\qquad
  g,\widetilde g\succeq\lambda h.
\end{equation}
Let $\delta_1=\|D\psi-D\widetilde\psi\|_{C^0_h}$ and
$\delta_2=\|\nabla^0D\psi-\nabla^0D\widetilde\psi\|_{C^0_h}$. Put
$A^\LC=\nabla^{g,\LC}-\nabla^0$ and
$\widetilde A^\LC=\nabla^{\widetilde g,\LC}-\nabla^0$. Then
\begin{align}
  \|g-\widetilde g\|_{C^0_h}
  &\le2B_1\delta_1,\\
  \|\nabla^0g-\nabla^0\widetilde g\|_{C^0_h}
  &\le2(B_1\delta_2+B_2\delta_1),\\
  \|A^\LC-\widetilde A^\LC\|_{C^0_h}
  &\le D_{\LC},
\end{align}
where
\begin{equation}
  D_{\LC}=\lambda^{-1}(B_1\delta_2+B_2\delta_1)
       +2B_1^2B_2\lambda^{-2}\delta_1.
\end{equation}
For every piecewise $C^1$ path,
\begin{equation}
  \|P_\gamma^{g,\LC}-P_\gamma^{\widetilde g,\LC}\|_{h\to h}
  \le L_\gamma D_{\LC}
  \min\left\{
    \exp\left(\frac{B_1B_2}{\lambda}L_\gamma\right),
    \frac{B_1^2}{\lambda}
  \right\}.
\end{equation}
If $\psi$ and $\widetilde\psi$ are $C^3$, their covariant Riemann tensors also
satisfy
\begin{equation}
  \|\operatorname{Rm}_g-\operatorname{Rm}_{\widetilde g}\|_{C^0_h}
  \le C(B_1,B_2,\lambda^{-1})(\delta_1+\delta_2),
\end{equation}
and the corresponding sectional-curvature conclusion of
Corollary~\ref{cor:sqrt-lc-curvature} holds under the same Gram-determinant
floor.
In particular the constants are independent of $\dim\mathbb H$.
\end{lemma}

\begin{proof}
The covariant Gram identities are
\begin{equation}
  g(X,Y)=\langle D\psi(X),D\psi(Y)\rangle,
  \qquad
  g(A^\LC(X,Y),Z)
  =\langle(\nabla^0D\psi)(X,Y),D\psi(Z)\rangle.
  \label{eq:covariant-square-root-identities}
\end{equation}
The second follows by inserting the first into the Koszul difference formula;
torsion-freeness makes $\nabla^0D\psi$ symmetric. Subtracting the two versions
of \eqref{eq:covariant-square-root-identities}, followed by the same
inverse-matrix identity as above, gives exactly the three bounds and
$D_{\LC}$ in Theorem~\ref{thm:sqrt-lc-stability}.

In an $h$-orthonormal $\nabla^0$-parallel frame along $\gamma$, the individual
coefficient norms are bounded by $B_1B_2/\lambda$, so Duhamel gives the
exponential branch. For the second branch,
$\lambda h\preceq g,\widetilde g\preceq B_1^2h$ and each Levi--Civita transport
is an isometry for its own metric. Hence each subpath propagator has
$h$-operator norm at most $B_1/\sqrt\lambda$; the exact variation identity
then gives $L_\gamma(B_1^2/\lambda)D_{\LC}$. Taking the smaller branch proves
the transport claim. For curvature, the tangent projection is
$\Pi=D\psi\,g^{-1}(D\psi)^*$ and the second fundamental form is
$(I-\Pi)\nabla^0D\psi$. The already proved inverse and first-derivative
estimates give
$\|\Pi-\widetilde\Pi\|\le C(B_1,\lambda^{-1})\delta_1$ and therefore control the
difference of the second fundamental forms by
$C(B_1,B_2,\lambda^{-1})(\delta_1+\delta_2)$. The Gauss equation in the flat
Hilbert target is quadratic in these forms, proving the covariant Riemann
bound. Choosing background-orthonormal plane representatives and using the
declared Fisher Gram floor proves the sectional-curvature statement.
\end{proof}

\subsection{The conditional-expectation projection}

\begin{lemma}[Orthogonal projection in the direct integral]\label{lem:direct-integral}
In the setting of Theorem~\ref{thm:semantic-obstruction}, let $\mathcal L$ be
the space of measurable sections $U$ with $U\in T_{F(H)}P$ almost surely and
$\E\|U\|_{g_{F(H)}}^2<\infty$, equipped with
$\langle U,V\rangle=\E[g_{F(H)}(U,V)]$, and let
$\mathcal L_F\subset\mathcal L$ be the subspace of sections of the form
$Y\circ F(H)$ with $Y$ measurable, where fields and sections are identified
almost everywhere under $\nu=\mathcal L(F(H))$. Then
$\E[\,\cdot\mid\sigma(F)]$ maps
$\mathcal L$ into $\mathcal L_F$ and is the orthogonal projection onto it; in
particular it is a contraction and $\mathcal L_F$ is closed.
\end{lemma}

\begin{proof}
Conditional expectations are taken in the finite-dimensional ambient zero-sum
space $T_0$. This is legitimate because
$g_p(u,u)=\sum_a u_a^2/p_a\ge\sum_a u_a^2$, so Fisher square integrability
implies ambient Euclidean square integrability. Since the embedded manifold
$P$ is a standard Borel space, the Doob--Dynkin lemma represents every
$\sigma(F)$-measurable ambient vector as a measurable function of $F$, up to
$\nu$-null sets. Conditionally on $F(H)=p$, the section takes values in the closed
linear subspace $T_pP\subset T_0$, and a conditional expectation of a random
vector supported in a closed convex set lies in that set almost surely; hence
$\E[U\mid\sigma(F)]\in T_{F(H)}P$ and belongs to $\mathcal L_F$. Idempotence is
the tower property. Moreover, conditional Jensen applied to the fixed
positive quadratic form $g_{F(H)}$ gives
\begin{equation}
  \|\E[U\mid\sigma(F)]\|_{g_{F(H)}}^2
  \le \E[\|U\|_{g_{F(H)}}^2\mid\sigma(F)],
\end{equation}
so the map is a contraction on $\mathcal L$. For self-adjointness, observe that $g_{F(H)}$ is
$\sigma(F)$-measurable, so for every $V\in\mathcal L_F$
\begin{equation}
\begin{split}
    \langle\E[U\mid\sigma(F)],V\rangle
    &=\E\bigl[g_{F(H)}\bigl(\E[U\mid\sigma(F)],V\bigr)\bigr]\\
    &=\E\bigl[\E\bigl[g_{F(H)}(U,V)\mid\sigma(F)\bigr]\bigr]
     =\langle U,V\rangle,
\end{split}
\end{equation}
because both $g_{F(H)}$ and $V$ may be moved inside the conditional
expectation. A bounded idempotent map whose range satisfies
$\langle\Pi U,V\rangle=\langle U,V\rangle$ for all $V$ in that range is the
orthogonal projection onto it, and the range of an orthogonal projection is
closed.
\end{proof}

\section*{Code and manuscript availability}

The manuscript source, model-free predictive-geometry implementation, CPU tests, and experimental protocols are available at
\begin{center}
    \url{https://github.com/guybass/probalistic-latent-spaces}.
\end{center}

\bibliographystyle{plainnat}
\bibliography{references}

\end{document}
